\documentclass[11pt,reqno,a4paper]{amsart}
\usepackage[T1]{fontenc}
\usepackage{amssymb,booktabs}
\usepackage[hidelinks]{hyperref}
\allowdisplaybreaks[2]
\raggedbottom
\setlength{\emergencystretch}{1.8em}

\newtheorem{theorem}{Theorem}[section]
\newtheorem{proposition}[theorem]{Proposition}
\newtheorem{corollary}[theorem]{Corollary}
\newtheorem{lemma}[theorem]{Lemma}
\theoremstyle{remark}
\newtheorem{remark}[theorem]{Remark}
\newtheorem*{theoremT}{Theorem T}

\newcommand{\Fac}{\operatorname{Fac}}
\newcommand{\e}{\varepsilon}
\newcommand{\Y}{\mathcal Y}
\newcommand{\tern}[1]{\mathtt{#1}}
\newcommand{\R}{\mathbb R}
\newcommand{\Q}{\mathbb Q}
\newcommand{\Z}{\mathbb Z}
\newcommand{\dist}[1]{\lVert#1\rVert}
\newcommand{\rev}{\mathrm{rev}}
\newcommand{\word}[1]{\mathtt{#1}}
\newcommand{\HH}{\mathcal H}

\title[Powers and palindromes under growing multipliers]
{Long powers and palindromes under growing integer multipliers}
\newsavebox{\affilbox}
\address{School of Artificial Intelligence, Shanghai Jiao Tong University, Shanghai, China}
\email{mqyqwerty@sjtu.edu.cn}
\date{}
\subjclass[2020]{11J68, 11J86, 11J87}
\keywords{Algebraic irrational numbers, integer-base expansions, squares, palindromes, integer multipliers, Schmidt Subspace Theorem}
\hypersetup{pdfauthor={Qiyuan Ma},pdftitle={Long powers and palindromes under growing integer multipliers},pdfdisplaydoctitle=true}

\begin{document}
\sbox{\affilbox}{%
  \begin{minipage}{0.92\textwidth}
    \centering\normalfont\normalsize
    School of Artificial Intelligence,\\
    Shanghai Jiao Tong University, Shanghai, China\\
    ORCID: \href{https://orcid.org/0009-0000-9532-0821}{0009-0000-9532-0821}\\
    \href{mailto:mqyqwerty@sjtu.edu.cn}{mqyqwerty@sjtu.edu.cn}
  \end{minipage}}
\author[Qiyuan Ma]{Qiyuan Ma\\[6pt]\protect\usebox{\affilbox}}

\begin{abstract}
Let $\alpha$ be a fixed real algebraic irrational. We study the occurrence of long palindromes and complete powers in the integer-base expansions of the multiples $a\alpha$, allowing the multiplier $a$ to grow with the target length. For every sufficiently large length $L$ there exist a multiplier $a_L=o(b^L)$ and a starting position $m_L=o(L)$ such that $a_L\alpha$ has a palindrome of length $L$ beginning after $m_L$ digits. We give exact asymptotic counts at selectable phases, including counts of true palindromes with a prescribed prefix, and we quantify the cost of transferring a complete $s$-th power back to a smaller multiplier: the least positive multiplier $a^*(h,s)$ that produces a complete $s$-th power of root length $h$ satisfies $\log_b a^*(h,s)/(sh)\to1$. On many scales, the original expansion of $\alpha$ itself has no two equal length-$h$ blocks with distinct starts among its first $T(h)$ digits, so any such power in the original digits must begin after that prefix. The stronger error term in the palindrome-counting theorem uses a $2026$ preprint interval form of the quantitative Subspace Theorem; a weaker bound depending only on the published theorem of Evertse and Ferretti is stated in parallel. These results concern growing multipliers and do not resolve Bugeaud's Problem~10.42 for the original expansion.
\end{abstract}
\maketitle

\section{Introduction}\label{sec:intro}

For an integer $b\ge2$ and a real number $x$, write $\{x\}=x-\lfloor x\rfloor$ and $\dist{x}=\min_{n\in\Z}|x-n|$. Let $\mathsf S_b(x)$ mean that the canonical base-$b$ expansion of the fractional part of $x$ contains factors $VV$ with $|V|$ unbounded. Let $\mathsf P_b(x)$ mean that it contains palindromes of unbounded length. Bugeaud's Problem~10.42 asks whether each property holds for every real algebraic irrational $x$ and every integer $b\ge2$ \cite[p.~220]{Bugeaud}. Adamczewski and Rampersad posed the corresponding questions in base two \cite{AR}.

The digit block starting after $i$ places and of length $d$ has integer value
\begin{equation}\label{eq:digitblock}
 D_{i,d}(x)=\lfloor b^{i+d}x\rfloor-b^d\lfloor b^i x\rfloor.
\end{equation}
Leading zeros are included. For a word $W$ let $[W]_b$ denote its integer value. Deleting any fixed prefix preserves both $\mathsf S_b$ and $\mathsf P_b$. A long palindrome can be trimmed symmetrically. If all unbounded squares began in a finite set of starts, one start would have unbounded roots. The second half of a sufficiently large square at that start contains each earlier given square at the same start, and moves that square beyond the deleted segment.

The long-repetition criterion of Adamczewski--Bugeaud--Luca, as stated by Bilu \cite[Theorem~3.3, p.~8]{Bilu}, implies that the greatest length of two disjoint equal factors among the first $N$ digits of a fixed real algebraic irrational is $o(N)$. The corresponding length $L_b(N)$ with distinct starts and possible overlap also satisfies
\begin{equation}\label{eq:known-early}
 L_b(N)=o(N).
\end{equation}
Indeed, overlapping matches of length $\ell$ at distance $d<\ell$ give a period-$d$ segment of length $\ell+d$. It contains $\lfloor\ell/d\rfloor+1$ adjacent copies of a length-$d$ word. Splitting them into two equal groups gives a disjoint match of length at least $\ell/3$, as in \cite[Lemma~3.4]{Bilu}.

The results of this paper concern changing positive integer multipliers of a fixed algebraic irrational, rather than the original expansion itself. In this broader setting one can produce long palindromes and complete powers, count them at selectable phases, and measure the cost of transferring a hit back to a smaller multiplier or to the original string.

Theorem~\ref{thm:pal-mult} produces, for every sufficiently large $L$, a multiplier $a_L=o(b^L)$ and a starting position $m_L=o(L)$ such that $a_L\alpha$ has a palindrome of length $L$ after position $m_L$. Subsequent sections give exact counts for growing templates, including complete powers at selectable phases and palindromes with a prescribed prefix, at a saving $b^{G(j)}$ over a full length-$L$ grid. Theorem~\ref{thm:late-separation} excludes repeated blocks of a prescribed length $h$ from the first $T(h)$ digits of the original expansion on many scales; a root of length $h$ in the original digits must therefore begin after that prefix. Theorem~\ref{thm:reduced-cost} and Corollary~\ref{cor:leading-cost} quantify the multiplicative cost of a complete $s$-th power of root length $h$: the least positive multiplier $a^*(h,s)$ satisfies
\[
 \frac{\log_b a^*(h,s)}{sh}\longrightarrow 1.
\]
Thus complete long powers can appear in some multiple $a\alpha$ with $a$ still smaller than the obvious full period $b^{sh}$, but a transfer of a complete $s$-th power back to the original expansion requires a multiplier whose size is essentially the length of the target block.

Theorem~\ref{thm:palindrome-counts} counts true palindromes, optionally with prescribed square prefixes. Its stronger discarded-scale bound uses the preprint interval theorem of Stephan \cite{Stephan}. The same statement is given with a weaker discarded-scale bound that depends only on the published theorem of Evertse and Ferretti \cite{EF}. For any use of the counting theorem that requires only a set of logarithmic density one of scales, the published bound is already sufficient whenever it is $o(J)$.

These statements are not the assertions of Problem~10.42. The original questions, for the expansion of $\alpha$ itself and for every integer base, remain open.

\section{Exact patterns under growing integer multipliers}\label{sec:positive-multipliers}

Let $e(t)=\exp(2\pi i t)$. A target set of integer labels $\mathcal A\subset\{0,\ldots,M-1\}$ represents the half-open intervals $[P/M,(P+1)/M)$ for $P\in\mathcal A$. Define
\[
 p=\frac{|\mathcal A|}{M},\qquad S(k)=\frac1{|\mathcal A|}
                      \sum_{P\in\mathcal A}e(-kP/M).
\]
We use a common smoothing calculation for both families of labels. For compactly supported, continuous, finitely piecewise linear functions $g,\phi$ on $\mathbb R$, let
\[
 F(x)=\sum_{P\in\mathcal A}\sum_{j\in\mathbb Z}\phi(M(x+j)-P),
 \qquad G(t)=\sum_{n\in\mathbb Z}g(n/M)e(nt).
\]
Then
\begin{align}
 \widehat F(k)&=p\,\widehat\phi(k/M)S(k),
 &|\widehat\phi(u)|&\ll_\phi(1+|u|)^{-2},\label{eq:sharedfourier}\\
 |G(t)|&\ll_g\min\{M,(M\|t\|^2)^{-1}\}.\label{eq:sharedg}
\end{align}
The first line follows by change of variables and integration by parts on the linear pieces. For the second line put $\delta=1/M$. The sum of the absolute second differences of $g(n\delta)$ is at most $\delta\operatorname{Var}(g')$, by integrating translated differences of $g'$ on intervals of length $\delta$. Taking the Fourier transform of the second difference gives the claimed decay away from the integers.

Assume that for all $0<|j|\le c_1 M$ one has $\|j\beta\|\ge c_2/M$, with positive constants $c_1,c_2$ independent of $M$. Split $|k|\le RM$ into $O(R)$ groups of length at most $c_1M$. In each group at most $O(1)$ of the distances $\|k\beta\|$ lie between $i/M$ and $(i+1)/M$. Summing \eqref{eq:sharedg} gives, for $R\ge1$ and integers $C\ge2$,
\begin{equation}\label{eq:shared-spacing}
 \sum_{|k|\le RM}|G(k\beta)|\ll RM,\qquad
 \sum_{\substack{|k|\le RM\\ \|k\beta\|>C/M}}|G(k\beta)|
                    \ll RM/(C-1).
\end{equation}
The constants may depend on $g,c_1,c_2$. The first contribution comes from $i=0$ and the convergent sum of $i^{-2}$ over $i\ge1$.

For a shift $\theta$, the weighted count is $\sum_k\widehat F(k)e(k\theta)G(k\beta)$. The zero frequency is
$pG(0)\int\phi$ with $G(0)/M\to\int g$. By dyadic summation of \eqref{eq:sharedfourier} and \eqref{eq:shared-spacing}, frequencies $|k|>RM$ contribute $O(pM/R)$. The frequencies $0<|k|\le RM$ with $|S(k)|<\rho$ contribute $O(\rho RpM)$. It remains to control the high coefficients $|S(k)|\ge\rho$ for fixed $R,\rho$. All estimates are independent of $\theta$. Once the remaining contribution is $o(pM)$, approximate indicators of the multiplier interval and of $[0,1)$ by nonnegative piecewise linear minorants and majorants. Their integrals can be as close as desired to the respective interval lengths, including pointwise inequalities at the half-open endpoints. This yields the unweighted count.

\subsection{Palindromes for arbitrary real algebraic irrationals}

\begin{theorem}\label{thm:pal-mult}
Let $x$ be a fixed real algebraic irrational and $b\ge2$. For every sufficiently large $L$ there are integers $m_L\ge0$, $a_L\ge1$ with $m_L=o(L)$ and $a_L=o(b^L)$ such that the length-$L$ block of $a_Lx$ starting after position $m_L$ is a palindrome. The block can have least period $L-o(L)$. For quadratic $x$ one may use $m_L=0$ for all large $L$. For every algebraic irrational one may use $m_L=0$ along an unbounded sequence of each parity.
\end{theorem}

\begin{proof}
Put $M=b^L$ and let $\mathcal A$ be the $b^{\lceil L/2\rceil}$ palindrome labels. The factorization of $S(k)$ over the mirror pairs contains the factors
\[
       b^{-1}\sum_{d=0}^{b-1}e(-dk w_i/M),
       \qquad w_i=b^i+b^{L-1-i},\quad 0\le i<\lfloor L/2\rfloor .
\]
There is also a central factor for odd $L$. Fix $\rho>0$ and take a small $\eta>0$ with $(b+1)\eta<b^{-2}$. Each index with $\|kw_i/M\|\ge\eta$ contributes a factor bounded in modulus by some $\gamma_{b,\eta}<1$. Thus $|S(k)|\ge\rho$ allows only $O(1)$ such indices. For two consecutive good indices use
\[
 bw_i-w_{i+1}=(b^2-1)b^{L-i-2},\qquad
 w_i-bw_{i+1}=-(b^2-1)b^i .
\]
Writing $E=(b^2-1)k\bmod M$, these give $\|E/b^{i+2}\|<b^{-2}$ and $\|E/b^{L-i}\|<b^{-2}$. In each case the corresponding two leading digits of the residue are both $0$ or both $b-1$. These conditions cover all neighboring pairs except for $O(1)$ pairs associated with bad indices or the center. The base-$b$ word of $E$ has $O(1)$ digit changes. Since $b^2-1$ is a unit modulo $M$, the number of such $k$ with $|k|\le RM$ is $O(L^J)$, for a fixed $J$ depending on $b,\rho,R$.

The same observation gives
\begin{equation}\label{eq:palresonance}
 \|k b^m x\|\ge b^{-\epsilon L}
 \quad(0<|k|\le RM,\ |S(k)|\ge\rho,\ 0\le m\le A_0 L)
\end{equation}
for any fixed $\epsilon,A_0>0$ and all sufficiently large $L$. Indeed, multiply the constant-digit-run representation for $E$ by $b-1$. For $|k|\le RM$, the integer $Dk$, with $D=(b-1)(b^2-1)$, is a bounded sum of integer multiples of powers of $b$, with bounded coefficients and largest exponent $O(L)$. In a counterexample sequence pass to a subsequence with fixed coefficients and constant or diverging adjacent exponent gaps. Merge the constant gaps and discard zero terms, so $Dk=\sum_i c_i b^{v_i}$ with consecutive gaps diverging. With $P$ the nearest integer to $k b^m x$, apply the $S$-adic Subspace Theorem to $(b^{m+v_i},P)$, using the archimedean form $(x/D)\sum_i c_i X_i-P$ in place of the last coordinate. Use the coordinate forms at primes dividing $b$. The power factors cancel and the form product is at most $b^{-\epsilon L}$. The height is at most $b^{O(L)}$. A fixed rational dependence on a subsequence would, after division by the largest power, involve the irrational limit $c_{\max}x/D$ of the scaled last coordinate. The two corresponding coefficients vanish. Divide by the remaining leading powers successively to see that all coefficients vanish. This proves \eqref{eq:palresonance}.

To obtain the spacing hypothesis, fix $Q=(b+2)M$. There is an $m=m_L=o(L)$ with a reduced rational $P/q$ satisfying
\begin{equation}\label{eq:reverse-dirichlet}
 Q/b<q\le Q,\qquad
 \big|\{b^m x\}-P/q\big|<1/(Qq).
\end{equation}
For completeness choose a late phase $J$ and use Dirichlet approximation with denominator at most $Q$, including an integer part so the rational shadows $b^Jx$. Divide by $b$ one step at a time to move back through the actual expansion. The reduced denominator is nondecreasing, grows by at most a factor $b$ per step, and its part $q_\perp$ coprime to $b$ stays fixed. At the first step (including step zero) where it exceeds $Q/b$, its size is at most $Q$, and the error is still less than $1/(Qq)$. If no such step occurs, the approximation at phase zero has error less than $b^{-J}/(Qq_J)$, where $q_J$ is the final-phase denominator. Ridout's theorem \cite{Ridout} gives for its initial-phase denominator $q_0\le Q$ an error at least $c_\delta q_0^{-1-\delta}(q_0)_\perp^{-1}$. Since $(q_0)_\perp\le q_J$, this forces $b^{-J}\gg Q^{-\delta}$. For any fixed $\delta>0$, starting at $J=\lceil 2\delta\log_b Q\rceil+C_\delta$ precludes this. The earliest admissible phase therefore satisfies $m_L=o(L)$. Finally, \eqref{eq:reverse-dirichlet} implies $\|j b^m x\|\ge1/(2q)$ for $0<|j|\le q/2$, giving \eqref{eq:shared-spacing}.

For the high coefficients choose $0<\epsilon<1/2$ in \eqref{eq:palresonance}. The contribution is at most $O(p L^J M^{-1+2\epsilon})=o(pM)$, by \eqref{eq:sharedg}. The shared smoothing calculation shows that for every fixed $c>0$ there are asymptotically $c|\mathcal A|$ positive $a<cM$ hitting the target at $m_L$. For $c$ small enough, the same spacing estimate gives $O(1)$ hits in each cylinder. Words of length $L$ with palindrome symmetry and a period $t$ number at most $b^{(t+2)/2}$ because reflection modulo $t$ has at most two fixed points. Roots of period at most $(1-\epsilon)L$ therefore account for $o(|\mathcal A|)$ labels. A diagonal choice over $c,\epsilon\downarrow0$ gives the first assertions.

A quadratic irrational is badly approximable, so for $m=0$ the spacing hypothesis holds at all $L$. For any irrational, an unbounded sequence of convergent denominators $q$ permits $M\le q<b^2M$ for an unbounded sequence of $L$ of either parity. Along it the error $<q^{-2}$ gives the same spacing estimate without shifting. The resonance argument applies at $m=0$ as well.
\end{proof}


\subsection{Exact repetition counts for quadratic irrationals}

\begin{lemma}\label{lem:pellresidue}
Let $\beta$ be a real quadratic irrational, $b\ge2$ an integer and $s\ge2$ fixed. Set $B=b^h$, $M=B^s$ and $D=B-1$. For every fixed $R,C,A>0$, all sufficiently large $h$ admit no integers satisfying
\[
 0<|k|\le RM,\qquad k\equiv r\pmod D,\quad |r|\le A,
 \qquad \|k\beta\|\le C/M.
\]
\end{lemma}

\begin{proof}
Take $k>0$ by changing signs. If such pairs exist for unbounded $h$, take $P$ to be the nearest integer to $k\beta$. Let $A_0X^2+A_1X+A_2$ be an integral minimal polynomial with discriminant $\Delta$. Choose $\kappa\in\{-1,1\}$ with $2A_0\beta+A_1=\kappa\sqrt\Delta$. For $\xi=2A_0P+A_1k+\kappa k\sqrt\Delta$, its conjugate is $2A_0(P-k\beta)$. Thus $0<|\xi\bar\xi|\ll_{\beta,R,C}1$ and $|\xi|\asymp k\to\infty$. Let $\varepsilon>1$ be a fixed norm-one unit in $\mathbb Z[\sqrt\Delta]$. Reducing the absolute value of $\xi$ modulo powers of $\varepsilon$ gives $\xi=z\varepsilon^j$, with $1\le|z|<\varepsilon$ and $|\bar z|=O(1)$. These bounds restrict the two integral coordinates of $z$ to finitely many values. Pass to a subsequence with $z,r$ fixed, $j\to\infty$, and $j=O(h)$.

Write $y=\varepsilon^j$. The congruence on $k$ implies $D\mid k-r$, where
\[
 2\kappa\sqrt\Delta\, y(k-r)=zy^2-2\kappa\sqrt\Delta\, r y-\bar z.
\]
Factor the polynomial as $z(y-\lambda_1)(y-\lambda_2)$ in a fixed number field. Both roots are nonzero. Choose a finite set of places $S$ containing the archimedean places and the places at which any of $b,z,2\sqrt\Delta,\lambda_i$ fails to be a unit. Write $v_\nu=-\log|\cdot|_\nu$ with product-formula normalization. Outside $S$, the divisibility gives
\[
 v_\nu(D)\le \sum_{i=1}^2
         \min\bigl(v_\nu(b^h-1),\,v_\nu(y/\lambda_i-1)\bigr).
\]
This is because $v_\nu(D)\le v_\nu(y/\lambda_1-1)+v_\nu(y/\lambda_2-1)$ and all the terms are nonnegative. The finitely many indices with $y=\lambda_i$ can be removed.

For each $i$, apply the $S$-unit common-divisor estimate of Corvaja and Zannier \cite[Proposition~2 of the preprint]{CZ} to $(X,Y)=(b^h,y/\lambda_i)$. Both logarithmic heights are $O(h)$. For every $\epsilon>0$, outside a finite set and finitely many one-dimensional subgroups $X^a=Y^d$, the sum of the local common-divisor heights is at most $\epsilon\max(\operatorname{ht}(X),\operatorname{ht}(Y))$, where $\operatorname{ht}$ is logarithmic height. Each fixed subgroup with $a\ne0$ can occur only for bounded $h$, by comparing valuations above a prime factor of $b$. The valuation of $Y$ is fixed there because $\varepsilon$ is a unit. If $a=0$, the archimedean growth of $Y$ excludes the subgroup eventually. Hence the sum outside $S$ is $o(h)$. At the fixed finite places in $S$ above primes not dividing $b$, the valuations of $b^h-1$ are $O(\log(h+1))$ by lifting the exponent, including the usual separate formula at $2$. At primes dividing $b$ the valuations are zero. Summing yields $\log D\le o(h)$, contrary to $\log D\sim h\log b$.
\end{proof}

\begin{theorem}\label{thm:quadcount}
Fix a real quadratic irrational $\alpha$, an integer $b\ge2$ and an integer $s\ge2$. Write $B=b^h$, $M=B^s$, and $L_s=(M-1)/(B-1)$. For fixed bounded half-open intervals $I\subset\mathbb R$ and $J\subset[0,1)$,
\begin{multline}\label{eq:repeatcount}
 \#\bigl\{n\in\mathbb Z:n/M\in I,\ \
     \lfloor M\{\theta+n\alpha\}\rfloor=L_s w,\
     0\le w<B,\ w/B\in J\bigr\}\\
                     = B|I||J|+o(B)\qquad (h\to\infty),
\end{multline}
uniformly in $\theta\in\mathbb R$. In particular for every large $h$ one can choose $a_h=o(B^s)$, a positive integer coprime to $b$, such that $a_h\alpha$ begins with a true $s$-fold repeat $V^s$ with $|V|=h$. The least period of $V$ can be $h-o(h)$ and the least period of the whole repeat is then exactly $h$. Any fixed starting position is allowed.
\end{theorem}

\begin{proof}
The case $|J|=0$ is empty. Otherwise let $\mathcal A$ consist of the labels $L_sw$ with $w/B\in J$. The root indices form one integer interval of length $B|J|+O(1)$. For fixed $\rho>0$, the geometric sum gives
\[
           |S(k)|\ge\rho\ \Longrightarrow\
             \|kL_s/M\|\ll_{J,\rho}B^{-1}.
\]
With $\ell$ the nearest integer to $kL_s/M$, take $r=k-\ell(B-1)$. The equality $M=(B-1)L_s+1$ gives $L_s r-\ell=kL_s-\ell M$. When $|k|\le RM$, $\ell$ has size at most $RL_s+1/2$. Since $BL_s\ge M$, $r$ is bounded in terms of $R,J,\rho$. By Lemma~\ref{lem:pellresidue}, for every fixed $C$ all these nonzero high coefficients eventually have $\|k\alpha\|>C/M$.

Bad approximability of $\alpha$ supplies the spacing hypothesis in \eqref{eq:shared-spacing}. Use its second estimate for the high coefficients. Their total contribution is $O(RpM/(C-1))$, with a constant independent of $C$. Let $C$ tend to infinity after fixing $R,\rho$, then use the low- and far-frequency bounds by taking $R$ large and $\rho$ small. This proves \eqref{eq:repeatcount} through the shared smoothing calculation.

For $0<a<cM$, restrict to $a\equiv1\pmod b$ by applying the count with step $b\alpha$ and shift $\alpha$. There are asymptotically $cB/b$ successes. The spacing of the points $\{a\alpha\}$ allows only $O_{\alpha,c}(1)$ of these in each interval of length $1/M$. There are $O_b(b^{(1-\epsilon)h})$ possible roots of least period at most $(1-\epsilon)h$. These roots therefore do not exhaust the successes. Make a diagonal choice over $c,\epsilon\downarrow0$. If $V^s$ has a period $t<h$, comparison over the first $h$ positions makes $V$ invariant under cyclic shift by $t$. Then $\gcd(t,h)\le h/2$ is a period of $V$, which contradicts the choice for large $h$. A fixed offset $m$ is accommodated by replacing $\alpha$ with $b^m\alpha$.
\end{proof}

A related rational obstruction gives an entire family of exact grids missing the target at one scale. For $1\le q<L_s$ set $Q=q(B-1)+1$. Scaling $[w L_s/M,(wL_s+1)/M)$ by $Q$ places it in $(wq,wq+1]$ for $0<w<B$, with the right endpoint of the scaled interval excluded. When $w=0$ the scaled interval contains only the integer $0$. Thus every nonzero $j/Q$, $0<j<Q$, misses the target. If $\gcd(Q,b)=1$, every shifted nonzero grid point misses the same target. These rationals depend on the scale.

For an original algebraic expansion an unbounded sequence of squares of root length $h$ must start after positions $m$ with $m/h\to\infty$, by \eqref{eq:known-early}. Setting $a_h=b^m$ in the bound $a_h=o(b^{2h})$ would give $m\le2h$, so Theorem~\ref{thm:quadcount} cannot simply be turned into a shift along the original expansion.

\section{Fixed-phase powers and fractional repetitions}
\label{sec:fixed-phase}

The conditions on the phase in Section~\ref{sec:phase-counts}
differ from those in the following two results. The first gives a
small power saving at a prescribed phase for every quadratic
irrational. The second gives prescribed-phase counts for a
shortened comparison in arbitrary degree.

\begin{theorem}\label{thm:fixed-quad}
Fix a real quadratic irrational \(\alpha\) and \(s\ge2\).
There is \(\eta>0\) such that, with \(A=B^{s-\eta}\), the
following assertion holds for fixed bounded \(I\) of positive
length. Let \(J_h\subset[0,1)\) be a half-open interval of length
at least \(B^{-\eta}\) and \(1\le q\le B^\eta\). For all real
\(\theta\) and
\[
 \beta=(u b^m\alpha+v)/d,\quad
 1\le |u|,d\le B^\eta,\quad 0\le m\le\eta h,
 \qquad u,v,d,m\in\Z,
\]
the number of integers \(n/A\in I\) for which \(\{\theta+n\beta\}\)
starts with \(V^s\), \([V]_b/B\in J_h\), and
\([V]_b\equiv w_0\bmod q\), is
\[
 (1+o(1))\, \frac{AB |I|\,|J_h|}{M q}.
\]
The error is uniform in the stated ranges. In particular at phase
zero there are powers of least period \(h\) with positive multipliers
\(a\equiv1\bmod b,\ a=o(B^{s-\eta})\). For squares, the root can be
prescribed to begin with \(U\) or \(cU\) and to end with
\(\overleftarrow U\), where \(U\) is any word of length
\(\lfloor\eta h\rfloor-1\). These choices give an even or an odd
palindrome at the seam, respectively.
\end{theorem}

\begin{proof}
We first show there exists \(\zeta>0\) such that
\[
 0<|k|\le B^{s+\zeta},\quad |r|\le B^\zeta,\quad
 k\equiv r\bmod D,\quad \|k\alpha\|\le B^{-s+\zeta}
\]
has no solution for large \(h\). Otherwise take a sequence of
counterexamples with \(\zeta\to0\), \(h\to\infty\), and \(k>0\).
The quadratic estimate \(\|k\alpha\|\gg 1/k\) gives
\(\log k=(s+o(1))\log B\).
Let \(a_0 X^2+a_1X+a_2\) be an integer minimal polynomial,
\(\Delta=a_1^2-4a_0a_2\), and \(2a_0\alpha+a_1=\sigma\sqrt\Delta\)
with \(\sigma=\pm1\). For the nearest integer \(P\) to \(k\alpha\),
put \(Z=2a_0P+a_1 k+\sigma k\sqrt\Delta\). Its conjugate is
\(\overline{Z}=2a_0(P-k\alpha)\), and
\(Z-\overline Z=2\sigma k\sqrt\Delta\).
Hence \(|Z|\asymp k\) and \(0<|Z\overline{Z}|\ll B^{2\zeta}\).
Use a fixed norm-one Pell unit \(\epsilon>1\) of
\(\Z[\sqrt\Delta]\) and write \(Z=z\epsilon^j\) with
\(1\le |z|<\epsilon\).
Then \(j\log\epsilon=(s+o(1))\log B\), and the coefficients \(z\),
\(\overline z\), and \(r\) have logarithmic height \(o(\log B)\).

Take \(S\) to consist of the archimedean places and those over
primes dividing \(b\) in \(\Q(\sqrt\Delta)\).
The two arguments \(X=B,\ Y=\epsilon^j\) are \(S\)-units of
maximum logarithmic height \(H\asymp\log B\).
The moving polynomials \(X-1\) and
\(zY^2-2\sigma\sqrt\Delta\,rY-\overline z\) are coprime of
fixed positive degrees with nonzero constant terms and height
\(o(H)\). The second value equals
\(2\sigma\sqrt\Delta\,Y(k-r)\), divisible by \(B-1\) at the
finite places in the ring of integers. Their logarithmic gcd
is therefore at least \(\log(B-1)\).
Grieve--Wang's moving gcd theorem
\cite[Theorem~1.2]{GrieveWang}, with sufficiently small
gcd tolerance, implies on an infinite subsequence that
\(B^a(\epsilon^j)^d\) has height \(o(H)\) for a fixed nonzero
integer pair \((a,d)\).
If \(a\ne0\), a valuation over a prime dividing \(b\) gives
height \(\gg h\), as the Pell unit has valuation zero.
If \(a=0\), the height is \(|d|j\operatorname{ht}(\epsilon)\gg h\).
Both cases contradict the conclusion. This proves the exclusion.

Choose \(0<\eta<\min(\zeta/32,1/16)\).
Let \(W=\{w:w/B\in J_h,\ 0\le w<B,\ w\equiv w_0\bmod q\}\).
It has cardinality \(B|J_h|/q+O(1)\gg B^{1-2\eta}\).
Smooth each digit cylinder at \(Lw/M\) on scale \(1/M\)
using a single compactly supported \(C^\infty\) function
\(\phi(Mx-Lw)\), periodized in \(x\). Their sum has coefficient
\[
 p\,\widehat\phi(k/M)S(k),\quad
 p=|W|/M,\qquad
 S(k)=|W|^{-1}\sum_{w\in W}e(-kLw/M).
\]
For the multiplier smoothing \(g(n/A)\), put
\(G_A(t)=\sum_n g(n/A)e(nt)\). Its magnitude is bounded by
\(O(A(1+A\|t\|)^{-8})\).
For \(1\le |k|\le T_1\) the quadratic estimate for \(\alpha\)
makes the sample points \(k\beta\bmod1\) separated by
\(\gg 1/(T_1 B^{3\eta})\).
Counting points in bands of width \(1/A\) gives, for \(T_1\ge A\),
\[
 \begin{aligned}
 \sum_{|k|\le T_1}|G_A(k\beta)|&\ll B^{3\eta}T_1,\\
 \sum_{\substack{|k|\le T_1\\ A\|k\beta\|>Z_1}}
       |G_A(k\beta)|&\ll B^{3\eta}T_1(1+Z_1)^{-7}
                  \qquad (Z_1\ge2).
 \end{aligned}
\]
For \(0<|k|\le K=MB^\eta\) with \(|S(k)|\ge B^{-8\eta}\),
the geometric sum gives
\(\|qkL/M\|\ll B^{-1+10\eta}\).
Taking the nearest integer \(l\) gives
\(r_1=qk-lD\) of magnitude \(O(B^{10\eta})\), using \(M=LD+1\).
If \(\|k\beta\|\le B^{-s+\zeta-2\eta}\), then the frequency
\(qk u b^m\) has magnitude at most \(B^{s+4\eta}\), residue
\(r_1 u b^m\) of magnitude \(O(B^{12\eta})\),
and its distance after multiplication by \(\alpha\) is
at most \(qd\|k\beta\|\le B^{-s+\zeta}\).
This is excluded by the first paragraph. Thus
\(A\|k\beta\|>B^{\zeta-3\eta}\) for these \(k\).
Relative to \(pA\), the error for \(|k|>K\) is
\(O(B^{3\eta}(M/A)(K/M)^{-7})=O(B^{-3\eta})\) by dyadic
summation and decay of \(\widehat\phi\).
For \(|k|\le K\) the small-\(|S(k)|\) and large-\(|S(k)|\)
errors are bounded by \(O(B^{-3\eta})\) and
\(O(B^{5\eta-7(\zeta-3\eta)})\), respectively.
The zero term yields \(pA\int g\int\phi+o(pA)\).
Sandwiching the half-open interval and the cylinders with
smooth functions proves the asserted exact count.

For the special multipliers take \(\beta=b\alpha,\theta=\alpha\).
Distinct tested multipliers have circle separation \(\gg 1/A>1/M\)
for large \(h\). There are \(\gg B^{1-3\eta}\) distinct roots
even with the prescribed templates, because the prefix interval
has width at least \(B^{-\eta}\) and the suffix prescribes one
residue with modulus at most \(B^\eta\).
Discard \(O(B^{1/2})\) roots with a period at most \(h/2\) to
obtain least period \(h\) for the full power. Allowing any fixed
positive proportion bound on \(a/A\) and taking a minimum yields
the little-oh assertion.
\end{proof}

\begin{theorem}\label{thm:fractional-fixed}
Fix any real algebraic irrational \(\alpha\), \(0<\xi<1\),
\(\delta,C>0\), and put \(r=\lfloor\xi h\rfloor\), \(E=b^r\),
\(A=E^2B^\delta\). For fixed bounded \(I\) of positive length,
among integers \(n/A\in I\) the number for which the first \(r\)
digits of \(\{\theta+n b^m\alpha\}\) agree with its length-\(r\)
block after \(h\) places is \((1+o(1))|I|A/E\), uniformly for
\(\theta\in\R\) and \(0\le m\le Ch\).
If \(0<\delta<1-\xi\) and
\(\xi+\delta>\max(1/2,1-\xi)\), a positive proportion of the hits
have least period \(h\) on the entire length-\((h+r)\) segment.
\end{theorem}

\begin{proof}
Let \(\kappa=\delta/16\), \(\lambda=\delta/4\).
For a fixed finite set of integers \(j\), all nonzero
\(k=jE+lD\), \(|l|\le EB^\kappa\), satisfy
\(\|k b^m\alpha\|>A^{-1}B^\lambda\) for large \(h\).
For \(l=0\) this is \eqref{eq:phase-ridout}, applied to
\(j b^{m+r}\). For \(l\ne0\) use the integer vector
\[
     (X_1,X_2,X_3,P)=(b^{m+h}l,b^m l,b^mE,P)
\]
where \(P\) is the nearest integer to \(k b^m\alpha\).
The real forms are the first three coordinates and
\(\alpha(X_1-X_2+jX_3)-P\). At primes dividing \(b\) use coordinate
forms. For a counterexample their product is at most
\(l^2 E^{-2}B^{-\delta+\lambda}\le B^{2\kappa-\delta+\lambda}\)
by cancellation of the base powers, and the height is \(B^{O(1)}\).
The \(S\)-adic Subspace Theorem \cite{Bilu} therefore puts an
infinite subsequence of counterexamples with fixed \(j\)
in a proper rational hyperplane. Divide its linear relation
by \(X_1\). The limit of \(P/X_1\) is \(\alpha\), while
\(X_2/X_1\) and \(X_3/X_1\) tend to zero. Irrationality
forces the coefficients of \(X_1,P\) to vanish.
The remaining relation gives \(l=c_0 E\) for a fixed nonzero
rational \(c_0\). Now \(k b^m\) equals
\(c_0 b^{m+r+h}+(j-c_0)b^{m+r}\).
Apply the same theorem on the vector
\((b^{m+r+h},b^{m+r},P)\) with error form
\(\alpha(c_0 X_1+(j-c_0)X_2)-P\).
The product is at most the error itself.
Dividing the new rational relation by its first coordinate
makes the limiting ratio of \(P\) equal to \(c_0\alpha\).
It forces all coefficients to vanish, a contradiction.

The exact block condition for \(x\in[0,1)\) is
\[
 -u\le t<1-u,\qquad u=\{Ex\},\quad t=E\langle Dx\rangle .
\]
Indeed \(\{Bx\}\) belongs to the same length-\(1/E\) cylinder
as \(x\) exactly when its signed circle displacement meets
these bounds. Use a trigonometric polynomial in \(u\), a smooth
test in \(t\), and a smooth multiplier test. The frequencies
are \(jE+lD\) and the nonzero terms with \(|l|\le EB^\kappa\)
have weighted error \(O(AB^{-N\lambda})\) by finite differences.
The tail is \(O(AB^{-\kappa(N-1)})\).
The zero term is \(A/E\) times the product of the test integrals
up to \(o(A/E)\). Taking \(N\) large and approximating the
region by rectangles as in Theorem~\ref{thm:selectable-counts}
yields the exact count, because its area is \(1\).

In the stated smaller range of \(\delta\),
\eqref{eq:phase-ridout} makes the hits use distinct
length-\((h+r)\) cylinders. For \(t_0<h\) a second period,
view a word of period \(t_0\) as labels modulo \(t_0\).
The equality constraints for period \(h\) join \(i\) to \(i+h\)
modulo \(t_0\) for \(0\le i<r\).
If \(r\ge t_0\), the graph has \(g=\gcd(h,t_0)\) components.
If \(r<t_0\), exactly \(\max(0,r-t_0+g)\) of the original shift
cycles remain complete, since their largest representatives are
\(t_0-g,\ldots,t_0-1\). There are then
\(t_0-r+\max(0,r-t_0+g)=\max(t_0-r,g)\) components.
Hence at most \(h b^{\max(h/2,h-1-r)}=o(A/E)\)
of all candidate words have a second period \(t_0<h\).
Discarding them proves the least-period assertion.
\end{proof}

Taking \(\xi=1-1/j\), \(\delta=1/(2j)\) and \(C=j\), and passing
through sufficiently large thresholds as \(j\to\infty\), gives
a comparison length \(r=h-o(h)\) with \(h-r\to\infty\).
The chosen positive multiplier is \(o(b^{h+r})\), and one
may take phase zero or a phase with \(m/h\to\infty\).
The comparison still leaves an intervening segment of \(h-r\)
digits.


\section{Exact power counts at selectable phases}\label{sec:phase-counts}

Fix an integer base \(b\ge2\) and a real algebraic irrational \(\alpha\).
For a root length \(h\) and integer \(s\ge2\) put
\[
 B=b^h,\qquad D=B-1,\qquad M=B^s,\qquad L=(M-1)/D.
\]
For an integer \(u>0\) let \(u_\perp\) denote its largest divisor coprime
to \(b\). Ridout's theorem gives, for every \(\epsilon>0\),
\begin{equation}\label{eq:phase-ridout}
 \|k b^m\alpha\|\gg_{\alpha,b,\epsilon}
               b^{-\epsilon m}|k|^{-1-\epsilon}
 \qquad (k\in\Z\setminus\{0\},\ m\ge0).
\end{equation}
Indeed the version with restricted denominator factors gives
\(q_\perp\|q\alpha\|\gg_{\alpha,b,\epsilon}q^{-\epsilon}\).
If the nearest numerator and \(q\) have a common factor, reducing the
fraction only decreases the product on the left
\cite{RidoutPadic}. Take \(q=|k|b^m\) and use \(q_\perp\le |k|\).

\subsection{Simultaneous exclusion of small-residue frequencies}

\begin{lemma}\label{lem:phase-residue}
Let \(s\ge2\) and \(C_0\ge0\) be fixed. Suppose
\(\kappa,\gamma_0,\rho>0\), \(\omega=\kappa+\gamma_0\),
\(2(\omega+\rho)<1\), and \(\rho<1/16\). For every sufficiently large
\(h\), every integer window
\([m_0,m_0+(2s+16)h]\), \(0\le m_0\le C_0h\), contains a run of
\(\gg h\) consecutive phases \(m\) at which the following system has
no solution:
\begin{equation}\label{eq:phase-bad}
 0<|k|\le B^{s+\kappa},\quad |r|\le B^\rho,\quad k\equiv r\pmod D,
 \qquad \|k b^m\alpha\|\le B^{-s+\gamma_0}.
\end{equation}
The excluded residues include \(r=0\).
\end{lemma}

\begin{proof}
Choose \(\sigma>0\) small and \(\Lambda\) so that
\[
 \sigma<1/16,\quad 2\omega+3\sigma<1,\quad
 \omega+\rho+3\sigma<\Lambda<1-\omega-\rho.
\]
Set \(G=\lfloor\sigma h/10\rfloor\), \(H=\lfloor\Lambda h\rfloor\),
\(d_1=2\omega+2\sigma\) and \(d_2=\omega+\rho+2\sigma\).
If there is no run of \(G\) good phases, choose one witness to
\eqref{eq:phase-bad} at each bad phase. Every interval of \(G\)
consecutive phases contains a witness. At relative offset \(j\) put
\[
 \alpha_0=b^{m_0}\alpha,\qquad Q_j=b^j k_j,\qquad
 R_j=b^j r_j,\qquad e_j=Q_j\alpha_0-P_j
\]
with \(k_j>0\) after a change of sign and \(P_j\) the nearest integer.
By \eqref{eq:phase-ridout}, with \(\epsilon\) small in terms of
\(C_0,s,\kappa,\sigma\), all the witnesses satisfy
\begin{equation}\label{eq:phase-range}
 k_j>B^{s-\gamma_0-\sigma},\qquad
 |e_j|\gg B^{-s-\kappa-\sigma/4}.
\end{equation}
The same bounds apply if the relative origin is reset at a subwindow.
In fact the only arithmetic inputs in the proof are
\eqref{eq:phase-range} and the lower bound on \(U+V\alpha_0\)
used below.

For two witnesses \(i<j\) with \(j-i\ge d_1h\), their
\((Q,P)\) projections cannot be proportional. Otherwise
\[
 |e_i|=(Q_i/Q_j)|e_j|
 \le B^{-(j-i)/h+\omega+\sigma}B^{-s+\gamma_0}
 \le B^{-s-\kappa-\sigma}.
\]
For nonzero residues, the full rows \((Q,P,R)\) cannot be proportional
if \(j-i\ge d_2h\). Proportionality would also imply
\(k_i/k_j=r_i/r_j\le B^\rho\), giving the stronger estimate
\(|e_i|\le B^{-s-\kappa-2\sigma}\).

For three witnesses spanning at most \(H\) positions let \(j_*\) be the
latest offset. The integer determinant of the rows
\((k_j,b^{j_*-j}P_j,r_j)\) is divisible by \(D\). Subtracting
\(b^{j_*}\alpha_0\) times the first column from the second bounds its
absolute value by \(6B^{\omega+\rho}b^H<D\) for large \(h\).
The determinant is therefore zero. Scaling rows and columns identifies
its rank with that of the original rows.

First consider a subwindow of length \((s+6)h+1\) with witnesses
throughout at the stated \(G\)-spacing, all of nonzero residue. Reset
its relative origin. Every \(H+1\) consecutive positions and every
overlap between two consecutive such intervals contains two witnesses
at distance at least \(d_2h\), by \(H-2G>d_2h\) with a fixed positive
margin. The rows consequently lie in a single rational plane
\(UQ+VP+WR=0\). Use a cross product from the first such window so that
\(|V|\le 2B^{s+\kappa+\rho+2\Lambda}\). The first and last witnesses
have independent \((Q,P)\) projections, so \(W\ne0\). We have
\begin{equation}\label{eq:plane-phase}
 (U+V\alpha_0)k_j+Wr_j=V b^{-j} e_j .
\end{equation}
If \(V\ne0\), the Ridout bound yields
\(\left|U+V\alpha_0\right|\gg B^{-\sigma/4}/|V|\)
by reducing \(\epsilon\) further. Using
\(2\kappa+2\rho+4\Lambda+\gamma_0<4\) gives
\[
 \frac{|V|B^{-s+\gamma_0}}{|U+V\alpha_0|}
    \ll B^{s+4+\sigma/4}.
\]
For \(j\ge(s+5)h\), all the ratios \(k_j/r_j\) are within
\(O(B^{-1+\sigma/4})=o(B^{-2\rho})\) of
\(-W/(U+V\alpha_0)\). Distinct ratios with denominators at most
\(B^\rho\) differ by at least \(B^{-2\rho}\).
Hence the ratios are equal, and the plane equation makes the
\((Q,P)\) projections proportional in this last segment. The segment
contains two witnesses at distance at least \(d_1h\), a contradiction.

If \(V=0\), the nonzero ratio \(r_j/k_j\) is fixed. Its reduced
denominator \(q\) divides each \(k_j\) and is greater than
\(B^{s-\gamma_0-\sigma-\rho}\). For adjacent witnesses of gap
\(g\le G\) the integer
\[
 k_iP_j-b^g k_jP_i=b^g k_j e_i-k_i e_j
\]
is divisible by \(q\) and has magnitude at most
\((b^G+1)B^\omega<q\). Hence it vanishes, giving proportional
projections throughout the subwindow, again a contradiction.

Return to the full window. For consecutive zero-residue witnesses
\(i<j\), if \(j-i\le H\) the same two-row determinant is divisible by
\(D\) and is smaller than \(D\). Their projections are proportional.
If \(j-i>H\), all intervening witnesses are nonzero. The intervals of
\(H\) interior positions and their adjacent overlaps transmit a single
plane by the same rank argument. The first and last intervals of
\(H\) interior positions join the respective endpoint row to that plane
by the three-row determinant. The plane intersects \(R=0\) in a line
because it contains a row with \(R\ne0\). Thus the projections of the
two endpoints are again proportional. All zero-residue witnesses now
have proportional projections. Their first-to-last span is less
than \(d_1h<h\). Removing that span leaves a subwindow of length
\((s+6)h+1\) at one end of the original interval, with only nonzero
witnesses and with the \(G\)-spacing. In the absence of zero
witnesses one can choose any such subwindow. The preceding
contradiction completes the proof.
\end{proof}

\subsection{Exact digit targets and root templates}

\begin{theorem}\label{thm:selectable-counts}
Fix \(s\ge2\), \(C_0\ge0\) and \(0<\gamma<1/2\). In every window
of Lemma~\ref{lem:phase-residue} there is a run of \(\gg h\)
consecutive phases on which the following count holds. Put
\(A=B^{s-\gamma}\). For bounded fixed half-open intervals
\(I\subset\R\), \(J\subset[0,1)\) of positive length and every
\(\theta\in\R\), count the integers \(a\) such that \(a/A\in I\)
and the first \(sh\) digits of \(\{\theta+a b^m\alpha\}\) are
\(V^s\), with \(w=[V]_b\) and \(w/B\in J\). The count is
\begin{equation}\label{eq:exact-power-count}
        (1+o(1))\,|I||J|\,AB/M .
\end{equation}
It is uniform in \(\theta\) and in the selected phases. Restriction
to \(a\equiv a_0\bmod d\), with \(d\) fixed, divides the main term
by \(d\).

There are also constants \(\chi,\tau>0\) such that one may select
a run of \(\gg h\) consecutive phases in the same window for the
following count. Put \(l=\lfloor\chi h\rfloor\), \(E=b^l\), and
\(E_1=E\) or \(bE\). Prescribe the first \(\log_b E_1\) digits and
the last \(l\) digits of \(V\). Suppose the center of the prescribed
prefix cylinder lies in \([1/4,3/4]\), and take any
\(1\le d\le B^\tau\), \(a_0\bmod d\). The count in that congruence
class with the prescribed templates and \(a/A\in I\) is
\begin{equation}\label{eq:exact-template-count}
        (1+o(1))\,\frac{|I|A}{d L E_1 E}.
\end{equation}
The error is uniform in the stated template and congruence ranges
and in \(\theta\) and the selected phases.
\end{theorem}

\begin{proof}
For a real \(y\) let \(\langle y\rangle\in[-1/2,1/2)\) denote its
centered fractional part. Since \(M=1+LD\), the exact criterion
at \(x\in[0,1)\), for large \(h\), is
\begin{equation}\label{eq:target-centred}
 \lfloor Mx\rfloor=Lw\text{ for some }0\le w<B
 \quad\Longleftrightarrow\quad -x\le L\langle Dx\rangle<1-x.
\end{equation}
A hit gives \(Dx-w=(Mx-Lw-x)/L\), of magnitude less than \(1/L\).
Conversely take \(w\) to be the nearest integer to \(Dx\).
At a hit \(w=\lfloor Bx\rfloor\).

For the first assertion take small positive \(\lambda,\kappa,\rho\)
with \(2(\kappa+\gamma+\lambda+\rho)<1\) and \(\rho<1/16\).
Use Lemma~\ref{lem:phase-residue} with \(\gamma_0=\gamma+\lambda\),
and write \(\beta=b^m\alpha\), \(e(t)=\exp(2\pi i t)\).
Take smooth compactly supported \(g,\psi\) and a fixed
trigonometric polynomial \(P\). For samples \(x_a=\{\theta+a\beta\}\)
sum
\[
 g(a/A)P(x_a)\sum_{j\in\Z}\psi(L(Dx_a+j)).
\]
The frequencies are \(u+vD\), with \(u\) from the finite support
of \(P\) and \(v\in\Z\). The periodic sum has coefficients
\(L^{-1}\widehat\psi(v/L)\). With \(a=a_0+dn\), the nonzero
frequencies \(|v|\le L B^{\kappa/2}\) give a step
\(d(u+vD)\) in the range of \eqref{eq:phase-bad}, and
\[
 (A/d)\|d(u+vD)\beta\|\gg_d B^\lambda.
\]
The weighted sum over the progression has bound
\[
 \left|\sum_n g((a_0+dn)/A)e(nt)\right|
 \ll_N (A/d)\bigl(1+(A/d)\|t\|\bigr)^{-N}
\]
by finite differences. The coefficient absolute sum is bounded.
The errors inside and outside the cutoff are, respectively,
\(O((A/d)B^{-N\lambda})\) and
\(O((A/d)B^{-(N-1)\kappa/2})\). Take a fixed \(N\) large enough
that both are \(o(A/(dL))\). The zero frequency has asymptotic
\((A/(dL))\int g\int_0^1P\int\psi\). None of these bounds depends
on \(\theta\). Trigonometric approximation and smooth
majorants and minorants for rectangles give Lebesgue area in
\((x,t)\), \(t=L\langle Dx\rangle\), after division by \(A/(dL)\).
A finite union of rectangles can enclose and approximate the domain
\(x\in J,\ -x\le t<1-x\) up to arbitrarily small area, since the
boundary has area zero. Enlarging or shrinking \(J\) by a fixed small
amount and then sending that amount to zero handles \(w/B=x+O(1/B)\).
The same approximation on \(I\) yields \eqref{eq:exact-power-count}
since \(1/L\sim B/M\).

For the second assertion choose \(\chi,\tau,\mu,\lambda>0\) with
\[
 \gamma+4\chi+3\tau+4\mu+\lambda<1/2,\qquad
 2\chi+\tau+2\mu<1/16 .
\]
Set \(\kappa=\chi+\tau+2\mu,\ \rho=2\chi+\tau+2\mu\),
\(\gamma_0=\gamma+\chi+\tau+\lambda\). These satisfy the hypotheses
of Lemma~\ref{lem:phase-residue}. Let \(x_0\) be the prefix cylinder
center and \(w_0\bmod E\) the specified suffix label. Use the smooth
product
\[
  P_h(x)T_h(x)=\varphi(E_1(x-x_0))
      \sum_{t\in\Z}\psi\bigl(L(Dx-w_0-Et)\bigr)
\]
on \(0\le x<1\), extended periodically in \(x\). The first factor
is supported away from both endpoints for large \(h\). On the
support of the second factor, the integer from its nonzero term
must be the nearest integer \(w\) to \(Dx\), with \(w\equiv w_0\bmod E\).

Expand the second factor in \(Dx\), with period \(E\), and write
its index \(\nu=Ev+j\) for \(0\le j<E\). The factor
\(P_h(x)e(-jx/E)\) has Fourier coefficients \(p_j(u)\) with
\[
 |p_j(u)|\ll_N E_1^{-1}(1+|u|/E_1)^{-N}
\]
uniformly, by integrating derivatives on its support of length
\(O(1/E_1)\). The full coefficient is
\((LE)^{-1}p_j(u)\widehat\psi((v+j/E)/L)e(-jw_0/E)\)
at the integer frequency
\[
 k=u+vD+(B/E)j,\qquad E k=(Ev+j)D+(Eu+j).
\]
On \(|u|\le E_1B^\mu,\ |v|\le LB^\mu\) with \(d\le B^\tau\),
the integer \(dEk\) has magnitude at most \(B^{s+\kappa}\) and
residue representative \(d(Eu+j)\) of magnitude at most \(B^\rho<D\).
Thus \(k=0\) there forces \(u=v=j=0\). If \(k\ne0\), the inequality
\(\|dk\beta\|\le B^{-s+\gamma+\tau+\lambda}\) would make \(dEk\)
a forbidden frequency. Hence the step satisfies
\((A/d)\|dk\beta\|>B^\lambda\). The full coefficient absolute sum
is \(O(1)\), since \((LE)^{-1}\sum_{j,v}
|\widehat\psi((v+j/E)/L)|=O(1)\) and \(\sum_u|p_j(u)|=O(1)\).
The tail outside the cutoff has absolute sum
\(O(B^{-\mu(N-1)})\). As above, arbitrary fixed derivative order
makes each error \(o(A/(dLE_1E))\). The zero term equals
\((1+o(1))\,A(\int g)(\int\varphi)(\int\psi)/(dLE_1E)\)
for nonzero integrals.

Finally the prefix condition is \(E_1(x-x_0)\in[-1/2,1/2)\)
and on its slight enlargement \(x=x_0+O(1/E_1)\).
One can sandwich this condition, the interval \(I\), and the
\(t\)-interval \([-x,1-x)\) in \eqref{eq:target-centred} between
smooth functions with arbitrarily small integral discrepancy.
For the \(t\)-interval use translates of smooth approximants to
\([0,1)\), with supports modestly enlarged and shrunken by a
fixed margin to accommodate \(x-x_0\). All derivative bounds used
above are uniform for \(x_0\in[1/4,3/4]\). This gives
\eqref{eq:exact-template-count} with the specified half-open digit
conditions.
\end{proof}

\begin{corollary}\label{cor:phase-palindrome}
For \(s=2\), \(\theta=0\) and the phases of
Theorem~\ref{thm:selectable-counts}, one can choose
\(a_h\equiv1\bmod b\), \(0<a_h=o(B^{2-\gamma})\), such that the
actual digits of \(a_h\alpha\) after \(m\) places contain a square
of least period \(h\). At its seam there is a prescribed
palindrome \(\overleftarrow{U}U\) or \(\overleftarrow{U}cU\) of
length \(2l\) or \(2l+1\), respectively, provided the initial
template \(U\) or \(cU\) has its cylinder center in \([1/4,3/4]\).
Here \(c\) is a digit, \(l=\lfloor\chi h\rfloor\), and
\(\overleftarrow{U}\) is the reversal of \(U\).
\end{corollary}

\begin{proof}
Prescribe the indicated prefix and the suffix \(\overleftarrow{U}\).
For any fixed \(0<c_1<1\), the number of hits in \(0<a<c_1A\)
with \(a\equiv1\bmod b\) is \(\gg_{c_1} B^{1-\gamma-2\chi}\).
Distinct multipliers in this range have circle distance greater
than \(1/M\) by \eqref{eq:phase-ridout} with sufficiently small
\(\epsilon\). Each root label specifies a digit cylinder of
length \(1/M\), so the hits use distinct roots. There are
\(O(B^{1/2})\) roots of length \(h\) with a period at most \(h/2\).
If \(V^2\) has a period \(t<h\), comparing positions in the two
copies makes \(V\) invariant under cyclic shift by \(t\). It then
has the period \(\gcd(t,h)\le h/2\). The parameter choice makes
\(1-\gamma-2\chi>1/2\), so there are hits of least period \(h\)
for every \(c_1\). The minimum successful positive multiplier
is \(o(A)\). The seam has the required digits because they are
part of the two completed copies.
\end{proof}


\section{Quantitative late phases}\label{sec:late-phases}

Let \(F:\Z_{\ge0}\to[1,\infty)\) be nondecreasing, with
\(F(J)\to\infty\) and
\[
 F(J)^2(1+\log F(J))^2=o(J).
\]
For \(h\ge1\) put \(j(h)=\lfloor\log_2h\rfloor\) and
\(T(h)=\lfloor h F(j(h))\rfloor\). An exceptional set of dyadic
indices will satisfy
\begin{equation}\label{eq:bin-count}
 \#\{j\le J:j\text{ is exceptional}\}
       \ll F(J)^2(1+\log F(J))^2 .
\end{equation}
A conclusion holds on a retained index simultaneously for all integer
\(h\in[2^j,2^{j+1})\). The corresponding retained integers have
logarithmic density one, since each complete dyadic interval has
harmonic weight \(\log 2+O(2^{-j})\).

\subsection{A uniform denominator bound and separation of blocks}

We use the common-weights interval theorem of Stephan
\cite[Theorem~4.2]{Stephan}, in the setting of fixed algebraic real
linear forms and coordinate forms at the primes dividing \(b\).
It extends the interval framework of Evertse--Ferretti
\cite{EF}. For a fixed system of \(n\) forms, suppose the exponents at the real
place are at most \(1\), with one equal to \(1\) and one at most
\(1-\delta/(2n)\). At each finite place they are nonpositive, with
one zero, and the sum of all exponents is at most \(-\delta\).
The inequality for each form is its absolute value bounded by the
sup norm of the integer vector raised to the corresponding exponent,
up to a fixed admissible real factor. Uniformly over the exponent
tuples satisfying these conditions, the theorem confines the
solutions with log height above \(O(\delta^{-1})\) to
\(O(1)\) proper rational subspaces or to log-height intervals
\([z,C\delta^{-1}z)\) of which there are
\(O(\delta^{-2}(1+\log\delta^{-1}))\).
All uses below have real factor \(1\) and real determinant of
absolute value \(1\). This quantitative input is a preprint result.

Suppose a counterexample at \(h\) has log height between
\(c2^{j(h)}\) and \(C(1+F(J))2^{j(h)}\), for \(j(h)\le J\),
with fixed \(c,C>0\). Use \(\delta=c_0/(1+F(J))\) for a
sufficiently small constant \(c_0>0\). A given log-height interval
allows a range of \(2^{j(h)}\) of ratio \(O((1+F(J))^2)\).
It meets \(O(1+\log F(J))\) possible indices \(j(h)\).
The height threshold accounts for \(O(1+\log F(J))\)
additional indices. It follows that the height intervals
contribute \eqref{eq:bin-count}. The points in proper subspaces
will be treated directly.

We will use the following elementary observation for the
three-dimensional cases. Write \(F_J=1+F(J)\). Consider points
\((b^\Delta,y)\) assigned to lengths \(h\) such that
\begin{equation}\label{eq:secant-hyp}
 c h\le\Delta\le C F_J h,\qquad
 c' b^{-C' F_J h}\le |y|\le b^{-c''h}
\end{equation}
for fixed positive constants. A single real affine line contains
such points with at most \(O(1+\log F(J))\) distinct dyadic indices.
Indeed, partition the indices into residue classes modulo
an integer \(O(1+\log F(J))\) so that lengths associated with
increasing indices in a class grow by a factor at least \(K F_J\),
with \(K\) sufficiently large. Except for a bounded number of small
lengths, this makes the horizontal coordinates grow by a factor
at least two and the absolute vertical coordinates fall by a factor
at least four. For three successive such points the horizontal
increment of the second secant is at least that of the first.
The first vertical increment has magnitude at least
\(3|y_1|/4\), while the second has magnitude at most \(5|y_1|/16\).
The secant slopes cannot agree. There are at most two indices
per class, which proves the observation.

\begin{theorem}\label{thm:late-separation}
For every fixed \(\xi>0\), outside a set of indices satisfying
\eqref{eq:bin-count}, all the retained \(h\) satisfy
\begin{equation}\label{eq:uniform-q}
       q_\perp\|q\alpha\|>B^{-\xi}
         \quad(1\le q\le b^{2T(h)}).
\end{equation}
In particular
\(\|k b^m\alpha\|>B^{-\xi}/|k|\) whenever
\(0\le m\le T(h)\) and \(0<|k|\le b^{T(h)}\).

With a possibly enlarged exceptional set satisfying
\eqref{eq:bin-count}, the first \(T(h)\) digits of the original
\(\alpha\) have no two equal \(h\)-digit factors at distinct starts,
with overlap allowed. In this prefix a square has root length less
than \(h\), and two palindromes of different centers have
intersection of length less than \(h\).
\end{theorem}

\begin{proof}
Use the fractional part of \(\alpha\), so \(0<\alpha<1\).
For a counterexample to \eqref{eq:uniform-q} let \(p\) be
the nearest integer to \(q\alpha\), and reduce \(p/q\) to
\(p'/q'\). Its error and its base-external factor have
both decreased, so \(q'_\perp |q'\alpha-p'|\le B^{-\xi}\).
The vector \((q',p')\) is primitive and has norm \(q'\).
Liouville's inequality gives
\(|q'\alpha-p'|\gg (q')^{-(\deg\alpha-1)}\), so the log height
is \(\gg h\), and it is at most \(2T(h)\log b\).
At the real place use \(X_1,\alpha X_1-X_2\) and exponents
\(1,\log|q'\alpha-p'|/\log q'\).
At \(p_0\mid b\) use exponents
\(\log|q'|_{p_0}/\log q',0\). Their total is
\(\log(q'_\perp|q'\alpha-p'|)/\log q'\), which is at most
\(-\delta\) for \(c_0>0\) small in terms of \(\xi\).
Each proper rational line contains at most one primitive
integer vector with positive first coordinate. Its height
can correspond to \(O(1+\log F(J))\) indices by the two height
bounds. The interval estimate then proves \eqref{eq:uniform-q}.
Use \(q=|k|b^m\) for its displayed consequence.

For the repeat assertion fix \(0<\xi<1/4\) in
\eqref{eq:uniform-q}. Let two length-\(h\) factors start after
\(m,n\) places, with \(m<n\) and \(n+h\le T(h)\).
Equality of their labels implies an integer \(p\) with
\[
 |e|=|\alpha(b^n-b^m)-p|<B^{-1}.
\]
If \(\Delta=n-m\le h/2\), the consequence of \eqref{eq:uniform-q}
with \(k=b^\Delta-1\) makes this impossible. Thus \(h/2<\Delta\le
T(h)\). The norm of \((b^n,b^m,p)\) is \(b^n\).
For the real forms
\(X_1,X_2,\alpha(X_1-X_2)-X_3\) use exponents
\(1,m/n,-h/n\). At \(p_0\mid b\) the exponents are
\(\log|b^n|_{p_0}/(n\log b),\,\log|b^m|_{p_0}/(n\log b),\,0\).
The total is \(-h/n\) and the log height
\(n\log b\) is between \(h(\log b)/2\) and \(T(h)\log b\).
Apply the interval estimate. A containing rational plane
\(uX_1+vX_2+wX_3=0\) with \(w=0\) fixes \(\Delta\), leaving
\(O(1+\log F(J))\) possible indices by \(h/2<\Delta\le T(h)\).
When \(w\ne0\), division by \(b^m\) puts
\((b^\Delta,e/b^m)\) on the line
\[
 y=(u/w+\alpha)x+(v/w-\alpha).
\]
Liouville's inequality with denominator \(b^n-b^m\) and the small
error bound give \eqref{eq:secant-hyp}, including constants
independent of the plane. Thus the number of exceptional indices
from the subspaces also satisfies \eqref{eq:bin-count}.

When a square of root length at least \(h\) is wholly in the
prefix, the first \(h\) digits of each copy are such a repeat.
If two palindromes of distinct centers intersect in at least \(h\)
digits, reflect a common \(h\)-digit factor about each center.
Its two reflections are equal factors at distinct starts in the
prefix. This proves the last assertions.
\end{proof}

For example, for any \(\nu>0\) one can take a nondecreasing
modification of
\[
 F(j)=\frac{\sqrt{1+j}}{(1+\log(1+j))^{1+\nu}}.
\]
Then \(T(h)\asymp h\sqrt{\log h}/(\log\log h)^{1+\nu}\), and
the bound \eqref{eq:bin-count} is \(O(J/(1+\log J)^{2\nu})\).
The absence of repeats is a scale-specific quantitative restriction
related to the fixed-positive-proportion exclusion of
Adamczewski--Bugeaud--Luca \cite[Theorem~3.3]{Bilu}.

\subsection{Selectable late phases and growing powers}

\begin{theorem}\label{thm:growing-power}
Fix \(0<\gamma<1/2\). Outside an exceptional set satisfying
\eqref{eq:bin-count}, the following assertions hold.
For fixed \(s\ge2\) and \(0<\theta_0<1\), the counts and seam
construction in Theorem~\ref{thm:selectable-counts} and
Corollary~\ref{cor:phase-palindrome} hold at runs of phases in
\[
 [\,\lfloor\theta_0T(h)\rfloor,\,
       \lfloor\theta_0T(h)\rfloor+(2s+16)h\,].
\]
One may also prescribe any integer \(s=s(h)\) in the range
\(2\le s\le F(j(h))/8\). There is a phase
\begin{equation}\label{eq:growing-phase-window}
 m\in[\,\lfloor T(h)/4\rfloor,\,\lfloor T(h)/4\rfloor+(2s+16)h\,]
\end{equation}
and a positive \(a\equiv1\bmod b\) with
\(a=o(B^{s-\gamma})\) for which \(a\alpha\) has an actual
\(V^s\) of least period \(h\) after \(m\) digits.
The little-oh is uniform for this range of \(s\).
\end{theorem}

\begin{proof}
Choose \(\xi>0\) small, and use the retained indices of
Theorem~\ref{thm:late-separation}. The proof of
Lemma~\ref{lem:phase-residue} uses the Ridout estimate for
\eqref{eq:phase-range} and the bound on \(|U+V\alpha_0|\).
For all the indicated phase windows \(m\le T(h)\). When
\(s\le F(j(h))/8\), the frequencies involved, including \(V\),
satisfy \(|k|,|V|\le b^{T(h)}\) for large \(h\).
The consequence of \eqref{eq:uniform-q} gives both inputs,
with \(\xi<\sigma/4\), and replaces the Ridout constants by uniform
bounds. The parameters \(G,H,\Lambda\) and the determinant
comparisons are independent of \(s\). The final segment
\(j\ge(s+5)h\) still suffices to compare the \(k_j/r_j\).
For fixed \(s\) the Fourier proof above applies directly.
Distinct multipliers of size \(O(A)\) have circle distance
greater than \(1/M\) by the same consequence of \eqref{eq:uniform-q},
using \(\xi<\gamma\).

For growing \(s\), we give the uniform Fourier estimates.
Let \(0<\varepsilon_0<1/100\) be fixed. Convolve a bounded
interval indicator with the \(p\)-fold convolution of the
uniform probability density on
\([-\varepsilon_0/p,\varepsilon_0/p]\).
The resulting smoothings have support in the fixed
\(\varepsilon_0\)-enlargement and total variation at most \(2\).
For derivative order \(k<p\) their \(L^1\) and \(L^\infty\)
derivative bounds are \(O((C p)^k)\). This follows by
differentiating \(k\) distinct factors of the convolution, each
becoming a signed measure of total variation \(p/\varepsilon_0\).
Their Fourier transforms therefore satisfy
\[
 |\widehat f(t)|\ll \min(1,(Cp/|t|)^k).
\]
Finite differences with sampling scale \(A'\) give
\[
 \left|\sum_n f((n+\zeta)/A')e(nt)\right|
   \ll A'\min(1,(Cp/(A'\|t\|))^k)
\]
uniformly in the sampling shift \(\zeta\) if \(k=o(A')\).
In fact the \(k\)-th difference has magnitude
\(O((Cp/A')^k)\) and only \(O(A'+k)\) sampled terms.
The bounded variation gives \(O(1)\) Riemann-sum error.

Fix \(c_1>0\), and choose \(g_p,\psi_p\) of this form
supported in \((0,c_1)\) and \((-1/12,1/12)\), respectively,
with fixed positive integrals. For \(g_p\) choose its smoothing
width small in terms of \(c_1\). Put \(P(x)=1/2-\cos(2\pi x)\).
If \(P(x)\sum_{j}\psi_p(L(Dx+j))>0\), then
\(1/6<x<5/6\) and \(|L\langle Dx\rangle|<1/12\).
Thus \eqref{eq:target-centred} gives a true hit.
Sum over \(a=bn+1\) with weight \(g_p(a/A)\).
At the chosen phases with \(\gamma_0=\gamma+\lambda\), the
nonzero frequencies \(u+vD\) with \(|u|\le1\) and
\(|v|\le LB^{\kappa/2}\) have
\((A/b)\|b(u+vD)\beta\|\gg B^\lambda\), as before.
The zero term is asymptotic to a positive constant times \(A/L\).
The sum \(L^{-1}\sum_v|\widehat\psi_p(v/L)|\) is \(O(p)\)
by the Fourier bound, splitting at a fixed multiple of \(pL\).
The internal error is \(O(Ap(Cp)^k B^{-k\lambda})\).
The tail is \(O(A(Cp)^k B^{-(k-1)\kappa/2})\),
by summing the power decay. Choose
\(k=\lceil C_1(s+1)\rceil\) with \(C_1\) sufficiently large, and
\(p=k+2\). Since \(\log p=o(h)\), both errors are \(o(A/L)\)
uniformly over the stated range of \(s\).

All weights are bounded, so there are \(\gg_{c_1} A/L\)
positive contributions. The hits use distinct root cylinders by
the circle distance bound just proved. Since \(A/L\asymp
B^{1-\gamma}\), one can discard every root with a period
at most \(h/2\), leaving powers of least period \(h\)
by the cyclic-shift argument of Corollary~\ref{cor:phase-palindrome}.
For every fixed \(c_1>0\) there is such a hit with \(0<a<c_1A\).
Taking the minimum proves the uniform little-oh statement.
\end{proof}

For \(s=\lfloor F(j(h))/8\rfloor\) the number of copies is unbounded
and \(sh\) is comparable to \(T(h)\). In \eqref{eq:growing-phase-window}
the complete block lies in the first \(T(h)\) digits for sufficiently
large \(h\). Theorem~\ref{thm:late-separation} excludes the same
root length in the original digits within that window on the retained
indices. The two results identify where a transfer from changing
multipliers needs additional information.


\section{Rational centers and the cost of growing powers}
\label{sec:rational-centers}

For \(x\in[0,1)\), a positive integer \(a\) is a hit at \((h,s)\)
if its fractional part \(\{ax\}\) starts with \(V^s\) for some
length-\(h\) word \(V\). Write \(w=[V]_b\), \(u=\lfloor ax\rfloor\),
and \(P=aD v=Du+w\). The corresponding rational center is
\(v=P/(aD)\). The digit cylinder for \(ax\) is equivalent to
\begin{equation}\label{eq:signed-centre}
        -\frac{w}{M}\le aD(x-v)<\frac{D-w}{M}.
\end{equation}
Indeed its fractional interval is
\([Lw/M,(Lw+1)/M)\), and \(LD=M-1\). The first inequality
is strict when \(x\) is irrational. In particular
\(|x-v|<1/(aM)\) for irrational \(x\).
The \emph{reduced cofactor} of a center is
\[
 q=\operatorname{den}(Dv)=\frac{a}{\gcd(a,P)},
\]
where \(\operatorname{den}\) means the positive reduced denominator.
It divides both \(a\) and \(\operatorname{den}(v)\).

\subsection{Uniform lower bounds for the reduced cofactor}

\begin{lemma}\label{lem:late-three-forms}
For each fixed \(\xi>0\), outside a set of dyadic indices
satisfying \eqref{eq:bin-count}, all the corresponding \(h\)
satisfy
\begin{equation}\label{eq:late-three-forms}
 (a_\perp)^2\,\|a(b^h-1)b^m\alpha\|>B^{-\xi}
 \quad (0\le m\le T(h),\ 1\le a\le b^{T(h)}).
\end{equation}
\end{lemma}

\begin{proof}
Normalize \(0<\alpha<1\). For a counterexample use the nearest
integer \(p\) to \(a(B-1)b^m\alpha\), and put
\(e=a(B-1)b^m\alpha-p\). Divide
\((Ba b^m,a b^m,p)\) by \(g=\gcd(a b^m,p)\) to get
\((B k,k,p')\) with \(k=a b^m/g\), \(p'=p/g\).
The sup norm is \(B k\), and the real forms are the three
used in Theorem~\ref{thm:late-separation}. Exponents may be
taken as the exact logarithmic ratios for the values \(Bk,k,|e|/g\)
at the real place, and for the first two coordinates at the primes
of \(b\), with zero for the third at those primes.
The sum, multiplied by \(\log(Bk)\), is
\[
 \log\bigl((k_\perp)^2 |e|/g\bigr)
     \le -\xi\log B,
\]
because \(k_\perp\le a_\perp\). The third real exponent is
nonpositive, and the first equals \(1\). The log height lies in
\([h\log b,(h+2T(h))\log b]\), permitting the same choice of
\(\delta\) and the interval argument.

For any rational plane containing these vectors, its equation
\(uX_1+v_0X_2+w_0X_3=0\) with \(w_0=0\) fixes \(B\).
With \(w_0\ne0\), the points
\((B,e/(a b^m))\) lie on the line
\[
 y=(u/w_0+\alpha)x+(v_0/w_0-\alpha).
\]
Here the upper bound in \eqref{eq:secant-hyp} follows from
\(|e|\le B^{-\xi}\), and the lower bound follows by Liouville's
inequality with denominator \(a(B-1)b^m\le b^{h+2T(h)}\)
and division by \(a b^m\). The horizontal exponent in
\eqref{eq:secant-hyp} is \(\Delta=h\). This gives
\(O(1+\log F(J))\) indices per plane. Combined with the
height-interval estimate, this proves the lemma.
\end{proof}

\begin{theorem}\label{thm:reduced-cost}
Fix \(\xi>0\) and use the retained indices of
Theorem~\ref{thm:late-separation} and Lemma~\ref{lem:late-three-forms}.
Let \(x=\{b^m\alpha\}\) with \(0\le m\le T(h)\) and suppose
\(a\) is a hit with \(aD\le b^{T(h)}\).
For \(q=\operatorname{den}(Dv)\) the following inequalities hold:
\begin{equation}\label{eq:reduced-cost}
 \begin{aligned}
 q q_\perp &> \frac{a M}{D^2} B^{-\xi}
                      \ \ge a B^{s-2-\xi},\\
 q q_\perp^2 &> \frac{a M}{D} B^{-\xi}
                      \ \ge a B^{s-1-\xi}.
 \end{aligned}
\end{equation}
In particular
\begin{equation}\label{eq:qperp}
 a_\perp\ge q_\perp>
      B^{\max\{s-2-\xi,\ (s-1-\xi)/2\}} .
\end{equation}
For \(0<\xi<1\) the eventual period length of the canonical
base-\(b\) expansion of \(v\) exceeds
\(h\max\{s-2-\xi,(s-1-\xi)/2\}\).
\end{theorem}

\begin{proof}
Since \(q Dv\) is integral, \eqref{eq:signed-centre} gives
\[
 \|qD b^m\alpha\|=\|qD x\|< (q/a)D/M .
\]
Apply \eqref{eq:uniform-q} to \(qD b^m\), which is at most
\(b^{2T(h)}\) and has base-external part \(q_\perp D\).
Apply \eqref{eq:late-three-forms} with \(a\) replaced by \(q\).
This gives both inequalities in \eqref{eq:reduced-cost}.
Using \(q\le a\) gives \eqref{eq:qperp}.
If \(R=\operatorname{den}(v)\), then \(q_\perp\mid R_\perp\).
After the base-prime part is removed by advancing past the
preperiod, the eventual period \(t\) is the multiplicative order
of \(b\) modulo \(R_\perp>1\). Thus \(b^t-1\ge R_\perp\ge
q_\perp\), as claimed.
\end{proof}

For example \eqref{eq:reduced-cost} implies the explicit cancellation
bound
\[
 \gcd(a,P)<\sqrt{a}\,D\,B^{\xi/2}/\sqrt{M}.
\]
When \(a\le B^{s-\gamma}\) this is at most
\(B^{1-\gamma/2+\xi/2}\), uniformly in \(s\).
For an elementary digit comparison at the rational center, assume
\(0<\xi<1\). Then \(R_\perp>1\), and every \(b\)-adic endpoint
\(j/b^n\) is at distance at least \(1/(R b^n)\) from \(v\).
Consequently \(|x-v|<1/(aM)\) ensures that \(x\) and \(v\)
share the first \(n\) digits whenever \(b^n\le aM/R\).
In particular this gives \((s-1)h\) common leading digits, as
\(R\le aD\). For hits with \(a\le B^{s-\gamma}\) the
upper endpoint \(n=\lfloor\log_b(aM/R)\rfloor\) of this
guarantee is at most \((s+2+\xi-\gamma)h\) by
\(R\ge R_\perp\ge q_\perp > B^{s-2-\xi}\).
For \(s>6+3\xi\) this endpoint is shorter than two eventual
periods of \(v\), even before accounting for the preperiod.
The guarantee comes only from the cylinder bound. Actual
agreement could extend further.

\begin{corollary}\label{cor:leading-cost}
On the retained indices, let \(s=s(h)\to\infty\) with
\(2\le s\le F(j(h))/8\). Let \(a_*(h,s)\) be the least
positive integer multiplier giving a complete \(s\)-fold power
with root length \(h\) at some phase \(0\le m\le T(h)\).
The same conclusion holds if the minimum is taken only over the
phase window \eqref{eq:growing-phase-window}, with
\(a\equiv1\bmod b\) and least period \(h\). In either case
\[
       \frac{\log_b a_*(h,s)}{s h}\longrightarrow 1 .
\]
\end{corollary}

\begin{proof}
Theorem~\ref{thm:growing-power} gives
\(a_*(h,s)\le B^{s-\gamma}\) for all sufficiently large retained
\(h\), for any fixed \(0<\gamma<1/2\).
This upper bound ensures \(a_*D\le b^{T(h)}\).
The lower bound \eqref{eq:qperp} gives
\(\log_b a_*(h,s)>(s-2-\xi)h\). Divide by \(sh\).
\end{proof}

\subsection{When two multipliers have one center}

The next statement is elementary. Its arithmetic consequences
use the retained indices above. The square case with coprime
multipliers gives the threshold \(a+c\le B+2\).

\begin{proposition}\label{prop:common-centre}
Let \(x\) be irrational. If \(a,c>0\) are hits at the same
\((h,s)\) and
\begin{equation}\label{eq:centre-threshold}
       \frac{a+c}{\gcd(a,c)}\le L+1,
\end{equation}
then their rational centers coincide. If \(a\) is a hit of
center \(v\) and \(d\mid a\) with \(\operatorname{den}(Dv)\mid d\),
then \(d\) is also a hit with center \(v\).
In particular the common-center case transfers to \(\gcd(a,c)\)
and to \(\operatorname{den}(Dv)\).
\end{proposition}

\begin{proof}
In case of distinct centers, exchange \(a,c\) so that
\(v_a<v_c\), and divide \(a,c\) by \(g=\gcd(a,c)\) to get
\(a',c'\). Set \(r=a'P_c-c'P_a\), a positive integer.
From \eqref{eq:signed-centre} applied at both centers,
\[
       rM<c'(D-w_a)+a'w_c .
\]
The two sides are congruent modulo \(D\), since
\(P_a\equiv w_a\), \(P_c\equiv w_c\bmod D\), and \(M\equiv1\bmod D\).
Their gap is at least \(D\). Thus \(M+D\le(a'+c')D\),
and \(M+D=(L+1)D+1\), proving the first statement.

For the second, put \(t=a/d\). If \(x>v\), choose the residue
\(w_d\in[0,D-1]\) of \(dDv\) and take \(w_a\) from the hit.
The positive representative \(D-w_a\in[1,D]\) is at most
\(t(D-w_d)\) since these two positive integers are congruent
modulo \(D\). Thus the upper inequality for \(a\) in
\eqref{eq:signed-centre} gives the upper inequality for \(d\).
The lower inequality is automatic. If \(x<v\), use the
representative \(w_d\in[1,D]\) and the corresponding positive
\(w_a\in[1,D]\). Then \(w_a\le t w_d\), giving the lower
inequality for \(d\), while the upper is automatic. In either
case with \(u_d=(dDv-w_d)/D\in\Z\), the signed inequalities place
\(dx\) in the cylinder with integer part \(u_d\) and digit
label \(Lw_d\). This proves the transfer.
\end{proof}

\begin{proposition}[Sharpness]\label{prop:sharp-centre}
The bound \(L+1\) in \eqref{eq:centre-threshold} for equality of
centers is sharp for every \(h\) and \(s\).
For each fixed \(b\) and \(s\) the same bound for transfer to
the greatest common divisor is sharp for arbitrarily large \(h\).
\end{proposition}

\begin{proof}
For equality of centers, take \(a=1,\ c=L+1\) and irrational
\(x\in(L/(cM),1/M)\). The two root labels are \(0,1\) and their
centers are distinct.

For the transfer assertion choose \(1<a\mid D-1\) with
\(\gcd(a,L+2)=1\) and let \(c=L+2-a\).
This choice is possible for all large \(h\) for even \(b\),
with \(a=2\), because \(D-1\) is even and \(L+2\) is odd.
For odd \(b\), take \(a=B-2\), and choose \(h\) to be a
multiple of the multiplicative orders of \(b\) modulo the
primes \(p\mid 2^s+1\) with \(p\nmid b\). Indeed
\(L+2\equiv 2^s+1\bmod a\), and at each such prime \(B\equiv1\).
At a prime dividing both \(b\) and \(2^s+1\) one has \(B\equiv0\);
neither case allows \(p\mid a\).

Set \(w_a=0\), \(w_c=D-(D-1)/a\), so \(a w_c=1+D(a-1)\).
Choose \(1\le u<a\) with \(cu\equiv-1\bmod a\),
and \(v_0=(cu-a+1)/a\in\Z\). For the potential hits take
\(P_a=Du\), \(P_c=Dv_0+w_c\), so \(aP_c-cP_a=1\).
The proposed centers satisfy \(v_a=u/a<v_c\).
The right endpoint of the \(a\)-cylinder is greater than the
left endpoint of the \(c\)-cylinder precisely when
\(M<cD+a w_c\), and the latter sum is \(M+D\).
The right endpoint of the \(c\)-cylinder is at least \(v_c\),
so it is greater than the left endpoint \(v_a\) of the
\(a\)-cylinder. Their interiors therefore intersect.
Take irrational \(x\) in the intersection. Its distance to
\(u/a\) is less than \(1/(aM)\). Since \(a\mid D-1\) and
\(\gcd(cu,a)=1\), it follows that
\[
  \|Dx\|\ge 1/a-D/(aM)>D/M,
\]
using \(M>(a+1)D\). Thus \(1=\gcd(a,c)\) is not a hit.
Multiplying both multipliers by any positive \(g\) and replacing
\(x\) by \(x/g\) gives the same threshold for a prescribed
greatest common divisor \(g\).
\end{proof}

\begin{corollary}\label{cor:centre-obstructions}
Under the assumptions of Theorem~\ref{thm:reduced-cost} with
\(0<\xi<1\), two hits sharing a center cannot have a greatest
common divisor composed only of primes dividing \(b\).
Moreover if \(A D\le b^{T(h)}\) and \(2A\le L+1\), the number
of hits with \(1\le a\le A\) at one phase is at most
\[
 \min\!\left\{
      (A B^{2-s+\xi})^{1/2},\,
      (A^2 B^{1-s+\xi})^{1/3}
      \right\}.
\]
If this set is nonempty it has a unique common center and its
smallest multiplier is \(\operatorname{den}(Dv)\).
\end{corollary}

\begin{proof}
The reduced cofactor \(q\) divides the greatest common divisor
and has \(q_\perp>1\) by \eqref{eq:qperp}.
For the count, \(2A\le L+1\) implies \eqref{eq:centre-threshold}
for every pair.
All hits are multiples of \(q\), and \(q\) itself is a hit by
Proposition~\ref{prop:common-centre}. If \(a_{\max}\) is the
largest hit, the number is at most \(a_{\max}/q\).
In \eqref{eq:reduced-cost} one may substitute \(a_{\max}\)
and use \(q_\perp\le q\), so
\(q^2>a_{\max}B^{s-2-\xi}\) and
\(q^3>a_{\max}B^{s-1-\xi}\). The result follows using
\(a_{\max}\le A\).
\end{proof}

There is also a two-phase restriction. Suppose \(a\) hits at
\(m\), \(c\) hits at \(m+\Delta\), and both use root length \(h\)
and exponent \(s\). At the earlier tail \(x\) these are hits
of \(a\) and \(b^\Delta c\). Suppose these two hits of \(x\)
share the same center, and \(aD\le b^{T(h)}\).
Then \(q\mid\gcd(a,b^\Delta c)\) and the proof of
Theorem~\ref{thm:reduced-cost}, using the stronger precision
\(|x-v|<b^{-\Delta}/(cM)\), gives
\[
 b^\Delta c M B^{-\xi}
   < \min\{D^2 q q_\perp,\ D q q_\perp^2\}.
\]
Thus at retained \(h\) even a later phase sharing an earlier
center needs a substantial base-external common factor.
These restrictions concern this specific transfer through a
rational center. The existence of sufficiently late powers or
palindromes in the original expansion remains open.


\section{Counting palindromes with smaller integer multipliers}
\label{sec:palindrome-counts}

The palindrome result below improves the qualitative \(o(b^L)\) multiplier bound of Theorem~\ref{thm:pal-mult}. It concerns the actual expansion of a changing integer multiple of a fixed number. Fix \(b\ge2\) and a real algebraic irrational \(\alpha\). Write \(j=\lfloor\log_2 L\rfloor\), let \(F,G\ge1\) be nondecreasing functions of \(j\) with \(G(j)\to\infty\), and put
\[
 \begin{gathered}
 V(j)=F(j)G(j),\qquad M=b^L,\qquad A=M b^{-G(j)},\\
 T=\lfloor LF(j)\rfloor,\qquad H=\lceil L/G(j)\rceil .
 \end{gathered}
\]
Let \(s_b=1+\omega(b)\), with \(\omega(b)\) the number of distinct prime factors of \(b\).

\begin{theorem}\label{thm:palindrome-counts}
With Stephan's preprint interval theorem one can discard at most
\(O_{\alpha,b}(V(J)^2(1+\log V(J))^2)\) dyadic indices \(j\le J\).
Using the published Evertse--Ferretti theorem one can instead discard
\(O_{\alpha,b}(V(J)^{2+2s_b}(1+\log V(J))^2)\).
Assume the chosen bound is \(o(J)\). For every sufficiently large \(L\) in a retained dyadic interval and every integer \(m_0\in[0,T-H]\), there is an integer \(m\in[m_0,m_0+H]\) with the following property.

For any word \(U\) of length \(r\le\lfloor L/5\rfloor\), any fixed bounded half-open interval \(I\) of positive length, any \(\theta\in\mathbb R\), and any congruence \(a\equiv a_0\bmod d\) with \(1\le d\le b^{G(j)}\), the number of integers \(a\) with \(a/A\in I\) for which the first \(L\) digits of \(\{\theta+a b^m\alpha\}\) form a palindrome starting with \(U\) is
\begin{equation}\label{eq:palindrome-count-main}
       (1+o(1))\frac{|I|A}{d}\frac{b^{\lceil L/2\rceil-r}}{M}.
\end{equation}
The error is uniform in the phase window, \(U,\theta,d,a_0\), for each fixed \(I\). The retained lengths have logarithmic density one.
\end{theorem}

\begin{proof}
We may replace \(\alpha\) by its fractional part. For the set \(\mathcal P_{L,U}\) of palindromic integer labels with \(L\) digits and prefix \(U\), its density among length-\(L\) labels is \(p_U=b^{\lceil L/2\rceil-r}/M\). Its normalized Fourier coefficient is
\[
 S_U(k)=|\mathcal P_{L,U}|^{-1}
                   \sum_{P\in\mathcal P_{L,U}}e(-kP/M),
 \qquad e(t)=e^{2\pi it}.
\]
As in the factorization for palindromic integer sums in \cite{BHS}, the mirror pairs with weights \(w_i=b^i+b^{L-1-i}\) are independent. Our prescribed prefix fixes the \(r\) exterior pairs. Thus
\[
 |S_U(k)|\le
   \prod_{i=r}^{\lfloor L/2\rfloor-1}
   \left|b^{-1}\sum_{\ell=0}^{b-1}e(\ell k w_i/M)\right|.
\]
If \(|S_U(k)|\ge b^{-8G(j)}\), there are at most \(C_bG(j)\) free indices with \(\|kw_i/M\|\ge \tau_b=1/[4(b+1)b^2]\). This follows from the strict maximum of the geometric sum at integer arguments. Therefore the interval of free indices from \(\lceil L/4\rceil\) to \(\lfloor3L/8\rfloor\) contains a block of \(\ell\gg_b L/G(j)\) consecutive indices with \(\|kw_i/M\|<\tau_b\). Using
\(w_i-bw_{i+1}=-(b^2-1)b^i\) gives \(\|(b^2-1)k b^i/M\|<b^{-2}\) for each adjacent pair. These overlapping two-digit windows are all zero pairs or all \(b-1\) pairs. With \(E=(b^2-1)|k|\), this gives
\begin{equation}\label{eq:pal-spectrum-split}
 E=b^v u+r_0,\qquad |r_0|\le b^{v-\ell},\qquad
                       5L/8-1\le v\le3L/4 .
\end{equation}
Indeed \(v=L-i_0\) where \(i_0\) is the first good index on that interval. Take \(r_0\) to be the residue modulo \(b^v\) for a zero block and the residue minus \(b^v\) for a \(b-1\) block. For \(b^{7L/8}\le|k|\le M b^{4G(j)}\) this gives \(u\ge1\), \(|r_0|/(b^v u)\le b^{-\ell}\), and \(|u r_0|\ll_b b^{L+4G(j)-\ell}\).

Write \(Q_\perp\) for the largest divisor of \(Q\) coprime to \(b\), \(A_d=A/d\), and set \(R=Z=b^{4G(j)},\ \lambda=b^{-8G(j)}\). We first exclude, on the stated number of dyadic indices, all failures of
\begin{align}
 Q_\perp\|Q\alpha\|&> b^{-L/G(j)}
                       &&(1\le Q\le b^{2T}),                     \label{eq:pal-denom}\\
 A_d\|dk b^m\alpha\|&>Z
       &&(0<|k|\le RM,\ |S_U(k)|\ge\lambda,\ 0\le m\le T).
                                                               \label{eq:pal-resonance}
\end{align}
The dimensions of the Subspace Theorem applications are two and three. At the real place use the fixed forms \((X_1,\alpha X_1-X_2)\) or \((X_1,X_2,\alpha(X_1+X_2)-X_3)\). At each \(p\mid b\) use coordinate forms. The determinants have absolute value one.

A failure in \eqref{eq:pal-denom}, after taking the nearest integer and reducing to a fraction \(p/q\), has \(q_\perp|q\alpha-p|\le b^{-L/G(j)}\). This is the product of the corresponding two-dimensional forms, or an upper bound on it at the finite places. Liouville's lower bound on the nonzero error gives \(\log q\gg_\alpha L/G(j)\), and \(\log q=O_b(LF(j))\). Each proper rational line has at most one primitive integer vector with positive first coordinate, so a fixed line contributes \(O_{\alpha,b}(1+\log V(J))\) indices \(j\le J\).

In \eqref{eq:pal-resonance} take positive \(k\) by symmetry and put
\[
 Q=(b^2-1)d k b^m,\qquad |e_0|=|Q\alpha-P|=\|Q\alpha\|
                                 \ll_b b^{-L+6G(j)}
\]
for a nearest integer \(P\). If \(k<b^{7L/8}\), then \(Q_\perp\ll_b b^{7L/8+G(j)}\). If \(k\ge b^{7L/8}\) and \(r_0=0\), then \(Q_\perp\le d u\ll_b b^{L+5G(j)-v}\) by \eqref{eq:pal-spectrum-split}. In either case \(Q_\perp|e_0|\le b^{-cL}\) for some \(c>0\) and large \(L\). The same two-dimensional argument applies, with upper log height \(O_b(LF(j))\).

Otherwise use
\[
                  X_1=b^{m+v}d u,\qquad
                  X_2=b^m d r_0,\qquad X_3=P.
\]
Its sup norm \(K\) is \(X_1\) for sufficiently large \(L\), by \(0<\alpha<1\) and \(|r_0|/(b^v u)\le b^{-\ell}\). Cancellation of the base factors at the finite places bounds the product of forms by
\[
               (d|u|)(d|r_0|)|e_0|
                           \ll_b b^{-\ell+12G(j)}
                           \le b^{-c'L/G(j)}.
\]
The log height has lower bound \(\gg_b L\) and upper bound \(O_b(LF(j))\). In the primitive-vector version the error is \(e_0/g\), where \(g\) is the common divisor, so its decay still gives the lower height bound by Liouville on \((X_1+X_2)/g\) and \(P/g\). The product of forms is multiplied by \(g_\perp^{-n}\le1\) in dimension \(n\).

Here is the interval count in detail. In each case the log heights at \(j\le J\) are between \(c_1 L/G(J)\) and \(C_1 LF(J)\). The product bounds allow a uniform parameter \(\delta=c_2/V(J)\) for sufficiently small \(c_2>0\). For Stephan \cite[Theorem~4.2]{Stephan}, at the real place use the actual logarithmic ratios to \(K\). The first is one and the error has a negative exponent. At the finite places take the actual coordinate log ratios for the first coordinate or first two coordinates and zero for the last one. They are nonpositive and the product bound gives a sum at most \(-\delta\). The fixed forms and these weight vectors satisfy the common-weights theorem. It gives a bounded number of proper rational subspaces and \(O(\delta^{-2}(1+\log\delta^{-1}))\) log-height intervals with endpoint ratios \(O(\delta^{-1})\). Log heights below the theorem's threshold \(O(\delta^{-1})\) cover \(O(1+\log V(J))\) dyadic indices. Each stated interval covers at most \(O(1+\log V(J))\) dyadic indices too, by the upper and lower height inequalities.

For Evertse--Ferretti \cite[Theorem~3.3]{EF}, use primitive vectors. Negating the normalized log-form values divided by \(\log K\) gives nonnegative weights of sum at least \(n+\delta\), since the vector is primitive at the finite places. There is at least one zero weight at each place, from \(K=X_1\) and the existence of a unit coordinate at each finite place. Truncate weights at \(n+1\), round them down with mesh \(\delta/(2n s_b)\), and scale the full vector down so its total is \(n+\delta/2\). The rounded total is at least \(n+\delta/2\) even when a weight needed truncation. This creates \(O(\delta^{-(n-1)s_b})\) fixed weight vectors. Taking their negatives gives nonpositive exponents of sum \(-n-\delta/2\), with the fixed determinant and coefficient fields. Apply the theorem to each, which gives the log-height interval count of order \(O(\delta^{-2}(1+\log\delta^{-1}))\) and a proper subspace. The same log-height comparison applies.

For the remaining subspaces in dimension three, use a plane \(a_1X_1+a_2X_2+a_3P=0\). Put \(x=X_2/X_1\), \(y=e_0/X_1\). They are nonzero. Uniformly for \(j\le J\), the absolute values of \(x,y,y/x\) are at least \(b^{-C LF(J)}\), using the nonzero integers and Liouville for \(e_0\). Their upper bounds are respectively \(b^{-cL/G(J)}, b^{-L/2}, b^{-L/2}\). If \(a_3=0\), the plane fixes \(x\). Otherwise \((x,y)\) lie on a line. Split the dyadic indices into residue classes modulo \(\lceil C'(1+\log V(J))\rceil\). With \(C'\) sufficiently large, the values at successively larger indices satisfy \(|x'|\le|x|/4,\ |y'|\le|y|/4,\ |y'/x'|\le|y/x|/16\). The slope of a secant then has absolute value between \(3|y/x|/5\) and \(5|y/x|/3\) using its earlier endpoint. Two successive secants cannot have the same slope. Thus at most two such indices in each class come from a fixed plane, with at most one when \(x\) is fixed. Together with the previous interval count this gives the two asserted exceptional-index bounds. The bad set is the set of actual failures, so the conclusion is simultaneous in the template and all phases and residue classes.

We now choose the phase. Put \(Q_0=(b+2)M\) and apply Dirichlet's approximation at \(m_0+H\) to \(\{b^{m_0+H}\alpha\}\), with reduced initial denominator \(q_0\le Q_0\) and error \(<1/(Q_0q_0)\). The reverse denominator step of Berend and Boshernitzan \cite{BBshift} sends \(p/q\) to \((p/q+e)/b\) where \(e\) is the actual removed digit. Reduction changes the proposed denominator \(bq\) by a divisor of \(b\). The step keeps the base-external part and reduces the error by \(b\). If the denominator first crosses \(Q_0/b\), the new one is at most \(Q_0\) and the error is \(<1/(Q_0q)\). If none crosses, at \(m_0\) the reduced denominator \(q_H\) satisfies
\[
 (q_H)_\perp\|q_H b^{m_0}\alpha\|
      <b^{-H}q_H/Q_0\le b^{-H}/b.
\]
The input integer \(q_H b^{m_0}\) is at most \(b^{2T}\) for large \(L\), contradicting \eqref{eq:pal-denom}. Hence at some selectable \(m\) the rational approximation to \(\beta=\{b^m\alpha\}\) has \(q\asymp_b M\) and error \(<1/(Q_0q)\le1/q^2\).
 
For the count, smooth the multiplier interval with a smooth compactly supported \(g\) and define
\[
 W(t)=\sum_{n\in\mathbb Z}g((a_0+dn)/A)e(nt).
\]
Finite differences give \(|W(t)|\ll_N A_d(1+A_d\|t\|)^{-N}\) for any fixed integer \(N\). In groups of \(k\)'s of span at most \(q/(2d)\), the \(dk\beta\) modulo one are \(1/(2q)\)-separated, since the rational approximation has error at most \(q^{-2}\). Grouping by distances to the nearest integer in bands of width \(1/A_d\) bounds each group's sum by \(O_N(q)\) and its sum restricted to \(A_d\|dk\beta\|>Z\) by \(O_N(q(1+Z)^{1-N})\). Thus over \(|k|\le K\), \(K\ge M\), the corresponding bounds are \(O_N(dK)\) and \(O_N(dK(1+Z)^{1-N})\).

Insert a smooth compactly supported bump \(\phi(Mx-P)\) in each target cylinder and periodize. The resulting detector has Fourier coefficients \(p_U\widehat\phi(k/M)S_U(k)\). Its sum at \(\theta+(a_0+dn)\beta\), weighted by \(g((a_0+dn)/A)\), has zero-frequency term
\[
                 p_U A_d\left(\int\phi\right)\left(\int g\right)+O(p_U).
\]
Use the preceding grouped bounds for the tail \(|k|>RM\), for the internal small spectrum \(|S_U(k)|<\lambda\), and for the large spectrum with \eqref{eq:pal-resonance}. After division by \(p_U A_d\), the three errors are bounded by
\[
   O_N(d^2b^{G(j)}R^{1-N}),\qquad
   O_N(\lambda d^2b^{G(j)}R),\qquad
   O_N(d^2b^{G(j)}R(1+Z)^{1-N}).
\]
They tend to zero uniformly for \(N\ge4\) fixed. Take smooth nonnegative upper and lower functions for \([0,1)\) and \(I\), with arbitrarily close integrals. The pointwise comparisons hold at the endpoints, and the \(\theta\) factors have absolute value one. This proves \eqref{eq:palindrome-count-main}. The set of retained lengths has logarithmic density one since every retained dyadic interval is retained in full.
\end{proof}

For \(I=(0,c]\), \(\theta=0\), and \(a\equiv1\bmod b\), the count gives a positive hit for every fixed \(c>0\) eventually, hence the least such positive hit is \(o(b^{L-G(j)})\). Without the prefix constraint, one can require the palindrome's least period to exceed \(L-\lceil6G(j)\rceil\). A length-\(L\) palindrome of period \(t\) is determined by its digits on the orbits of \(i\mapsto L-1-i\bmod t\), at most \(t/2+1\) choices of positions. Summing \(b^{t/2+1}\) over the excluded periods gives \(O_b(b^{L/2-3G(j)})\) labels. The separation of multiplier outputs gives \(O_b(1)\) hits per label in a fixed \(I\), since its integer span is \(o(q)\). This discarded count is \(o(b^{L/2-2G(j)})\), below the main term for the empty prefix. With a prescribed prefix one can instead choose a square \(U=VV\) with \(|V|=\lfloor L/10\rfloor\) and \(VV\) of least period \(|V|\), giving both patterns in one multiple.

In the preprint version take \(F=1\) and \(G(j)\asymp\sqrt j/(\log j)^{1+\zeta}\), for any fixed \(\zeta>0\), with \(m_0=0\). This gives \(m=o(L)\) and the saving \(\exp(-c\sqrt{\log L}/(\log\log L)^{1+\zeta})\) relative to \(b^L\). Alternatively \(F(j)=(1+j)^\tau,\ G(j)=(1+j)^\gamma\) with \(\tau,\gamma>0\) and \(\tau+\gamma<1/2\) allow \(m/L\to\infty\) by taking \(m_0=\lfloor T/2\rfloor\). The published version allows \(\tau+\gamma<1/(2+2s_b)\) and preserves the same count with its slower admissible growth.

For quadratic \(\alpha\), bounded partial quotients supply \(q\asymp_\alpha M\) and the required approximant directly at \(m=0\) for every retained \(L\). In general the logarithms of consecutive convergent denominators have ratio tending to one by Roth's theorem \cite{RothClassic}. In each sufficiently large retained dyadic interval a convergent with \(\log_b q\) between \((5/4)2^j\) and \((7/4)2^j\) thus supplies both an even and an odd integer \(L\) immediately below \(\log_b q\), with \(q/b^L\in[1,b^2)\). In either case the arithmetic estimate \eqref{eq:pal-resonance} applies also at \(m=0\), so the Fourier argument gives the same count there. The count uses smoothing over changing multipliers. Specializing to the original digits would require a different average, such as estimates on \(e(k b^m\alpha)\) along the phase orbit at the scale of the palindromic target.


\section*{Numbering correspondence}
For ease of cross-reference, the following table records the numbering of the principal statements as they appeared in an unpublished extended account of related results dated 26~November~2026.
\begin{center}
\begin{tabular}{@{}ll@{}}
\toprule
This paper & Extended account \\
\midrule
Theorem~\ref{thm:pal-mult} & Theorem~7.1 \\
Lemma~\ref{lem:pellresidue} & Lemma~7.2 \\
Theorem~\ref{thm:quadcount} & Theorem~7.3 \\
Theorem~\ref{thm:fixed-quad} & Theorem~8.1 \\
Theorem~\ref{thm:fractional-fixed} & Theorem~8.2 \\
Lemma~\ref{lem:phase-residue} & Lemma~9.1 \\
Theorem~\ref{thm:selectable-counts} & Theorem~9.2 \\
Corollary~\ref{cor:phase-palindrome} & Corollary~9.3 \\
Theorem~\ref{thm:late-separation} & Theorem~10.1 \\
Theorem~\ref{thm:growing-power} & Theorem~10.2 \\
Lemma~\ref{lem:late-three-forms} & Lemma~11.1 \\
Theorem~\ref{thm:reduced-cost} & Theorem~11.2 \\
Corollary~\ref{cor:leading-cost} & Corollary~11.3 \\
Proposition~\ref{prop:common-centre} & Proposition~11.4 \\
Proposition~\ref{prop:sharp-centre} & Proposition~11.5 \\
Corollary~\ref{cor:centre-obstructions} & Corollary~11.6 \\
Theorem~\ref{thm:palindrome-counts} & Theorem~12.1 \\
\bottomrule
\end{tabular}
\end{center}

\begin{thebibliography}{99}
\bibitem{AR}
B. Adamczewski and N. Rampersad,
\emph{On patterns occurring in binary algebraic numbers},
Proc. Amer. Math. Soc. 136 (2008), 3105--3109.
\href{https://doi.org/10.1090/S0002-9939-08-09319-2}{DOI}.
\href{https://adamczewski.perso.math.cnrs.fr/Overlap.pdf}{Author version}.

\bibitem{Bugeaud}
Y. Bugeaud,
\emph{Distribution Modulo One and Diophantine Approximation},
Cambridge Tracts in Mathematics, vol.~193, Cambridge University Press,
2012.
\href{https://doi.org/10.1017/CBO9781139017732}{DOI}.

\bibitem{Bilu}
Y. F. Bilu, \emph{The many faces of the subspace theorem [after
Adamczewski, Bugeaud, Corvaja, Zannier...]}, S\'eminaire Bourbaki,
expos\'e 967, Ast\'erisque 317 (2008), 1--38.
\href{https://www.numdam.org/book-part/AST_2008__317__1_0/}{Publication record}.
\href{https://arxiv.org/pdf/0907.2098#page=8}{Theorem 3.3, printed p.~8}.

\bibitem{EF}
J.-H. Evertse and R. G. Ferretti, \emph{A further improvement of the Quantitative Subspace Theorem}, Ann. of Math. 177 (2013), 513--590. \href{https://doi.org/10.4007/annals.2013.177.2.4}{DOI}.

\bibitem{CZ}
P. Corvaja and U. Zannier,
\emph{A lower bound for the height of a rational function at $S$-unit points},
Monatsh. Math. 144 (2005), 203--224.
\href{https://arxiv.org/abs/math/0311030}{Author version}.

\bibitem{Ridout}
D. Ridout,
\emph{Rational approximations to algebraic numbers},
Mathematika 4 (1957), 125--131.
\href{https://doi.org/10.1112/S0025579300001182}{DOI}.

\bibitem{Stephan}
R. Stephan,
\emph{The Evertse--Ferretti interval result for a set of weights, and the exponent $1/2$ in three complexity bounds},
preprint (2026).
\href{https://arxiv.org/abs/2611.16536v1}{arXiv:2611.16536v1}.

\bibitem{GrieveWang}
N. Grieve and J. T.-Y. Wang,
\emph{Greatest common divisors with moving targets and consequences
for linear recurrence sequences},
Trans. Amer. Math. Soc. 373 (2020), 8095--8126.
\href{https://doi.org/10.1090/tran/8220}{DOI}.
\href{https://arxiv.org/abs/1902.09109}{Author version}.

\bibitem{RidoutPadic}
D. Ridout,
\emph{The \(p\)-adic generalization of the Thue--Siegel--Roth theorem},
Mathematika 5 (1958), 40--48.
\href{https://doi.org/10.1112/S0025579300001339}{DOI}.

\bibitem{BHS}
W. D. Banks, D. N. Hart, and M. Sakata,
\emph{Almost all palindromes are composite},
Math. Res. Lett. 11 (2004), 853--868.
\href{https://doi.org/10.4310/MRL.2004.v11.n6.a10}{DOI}.

\bibitem{BBshift}
D. Berend and M. Boshernitzan,
\emph{On a result of Mahler on the decimal expansions of \((n\alpha)\)},
Acta Arith. 66 (1994), 315--322.
\href{https://doi.org/10.4064/aa-66-4-315-322}{DOI}.
\href{https://matwbn.icm.edu.pl/ksiazki/aa/aa66/aa6642.pdf}{Journal text}.

\bibitem{RothClassic}
K. F. Roth, \emph{Rational approximations to algebraic numbers},
Mathematika 2 (1955), 1--20, corrigendum 168.
\href{https://doi.org/10.1112/S0025579300000644}{DOI}.


\end{thebibliography}

\end{document}
